\chapter{行为引导量化框架设计}

\section{注意力近似误差的代数分解与问题形式化}
\label{sec:ch3-problem}

在长上下文推理中，KV Cache 的量化误差不会被下游任务直接读取，而是先进入注意力打分、softmax 与 Value 聚合。本文因此不把 $K$ 与 $V$ 的逐元素重建误差作为唯一对象，而把量化前后的注意力分布与注意力输出一起作为需要保持的行为。本节只做两步：先固定单头 Decode 注意力的记号，再把输出误差拆成分布侧与聚合侧两项。

设当前解码步的 Query 向量为 $q \in \mathbb{R}^{d_k}$，历史上下文长度为 $S$，对应的 Key 与 Value 分别记为 $K=[k_1,\dots,k_S]^\top$、$V=[v_1,\dots,v_S]^\top$。标准单头注意力可写为
\begin{equation}
a=\mathrm{softmax}\!\left(\frac{qK^\top}{\sqrt{d_k}}\right), \qquad o=aV,
\end{equation}
其中 $a\in\mathbb{R}^{S}$ 表示注意力分布，$o$ 表示注意力输出。基于此，本文将注意力行为形式化为如下联合对象：
\begin{equation}
\mathcal{B}(q,K,V):=(a,o).
\end{equation}
在不引起歧义时，本文暂略层索引、头索引与时间步索引，以保持记号简洁；多层、多头情形可视为同一分析在不同层与不同头上的重复与聚合。

设 $Q_{\theta_K}(\cdot)$ 与 $Q_{\theta_V}(\cdot)$ 分别表示作用于 Key 与 Value 的量化映射，其参数由 $\theta_K$ 与 $\theta_V$ 表示。量化后的 Key 与 Value 记为
\begin{equation}
\hat K = Q_{\theta_K}(K), \qquad \hat V = Q_{\theta_V}(V),
\end{equation}
对应的量化后注意力分布与输出为
\begin{equation}
\hat a=\mathrm{softmax}\!\left(\frac{q\hat K^\top}{\sqrt{d_k}}\right), \qquad \hat o=\hat a \hat V.
\end{equation}
于是，KV Cache 量化在概念层上的目标可写为
\begin{equation}
\min_{\theta_K,\theta_V}\Delta_{\mathrm{beh}}\big((a,o),(\hat a,\hat o)\big),
\end{equation}
其中 $\Delta_{\mathrm{beh}}$ 表示量化前后注意力行为的联合偏移度量。式中对象是概念层目标，不要求后续每条路径共享同一个标量损失；实际执行时，对称路径使用注意力分布 KL，角色感知低比特路径则拆成 K-path 与 V-path 两个代理。

从量化输出与参考输出的定义出发，输出误差可作如下精确的代数展开：
\begin{align}
\hat o-o
&=\sum_{i=1}^{S}\hat a_i\hat v_i-\sum_{i=1}^{S}a_i v_i \\
&=\sum_{i=1}^{S}(\hat a_i\hat v_i-\hat a_i v_i)+\sum_{i=1}^{S}(\hat a_i v_i-a_i v_i) \\
&=\sum_{i=1}^{S}\hat a_i(\hat v_i-v_i)+\sum_{i=1}^{S}(\hat a_i-a_i)v_i.
\label{eq:ch3-error-decomp}
\end{align}
式~\eqref{eq:ch3-error-decomp} 是在第二行加减同一中间项得到的恒等式，不是把 Key 与 Value 的影响假设为相互独立。右侧第一项描述 Value 侧内容被量化后如何进入聚合结果；右侧第二项描述 Key 侧扰动经 logits 和 softmax 改变注意力分布后如何传到输出。

这个写法同时给出一个边界：Value 项中的权重是量化后的 $\hat a_i$，不是参考路径的 $a_i$。因此，Value 误差并不是在固定注意力分布下单独发生；如果 Key 量化已经改变了分布，聚合侧误差也会继承这种变化。后续使用分布侧代理和输出侧代理时，本文只把它们当作可执行的近似读数，而不把两条传播路径解释为完全隔离的因果分量。

% 图 3-1 建议置于此处
% 图题：图 3-1 注意力误差的两条耦合传播路径及其精确代数分解
\begin{figure}[!htbp]
  \centering
  \begin{tikzpicture}[
    x=1cm, y=1cm,
    >=Stealth,
    every node/.style={font=\small, align=center},
    arr/.style={-stealth, line width=0.95pt},
    qarr/.style={-stealth, line width=0.95pt, red!70!black, dashed},
    fp/.style={
      rectangle, rounded corners=3.5pt,
      draw=blue!58!black, line width=0.9pt,
      fill=blue!5, minimum width=2.75cm, minimum height=0.82cm
    },
    quant/.style={
      rectangle, rounded corners=3.5pt,
      draw=red!62!black, line width=0.95pt,
      fill=red!5, minimum width=2.35cm, minimum height=0.82cm
    },
    kpath/.style={
      rectangle, rounded corners=4pt,
      draw=cyan!55!black, line width=1.0pt,
      fill=cyan!8, minimum width=4.15cm, minimum height=1.06cm
    },
    vpath/.style={
      rectangle, rounded corners=4pt,
      draw=orange!78!black, line width=1.0pt,
      fill=orange!10, minimum width=4.15cm, minimum height=1.06cm
    },
    eqcard/.style={
      rectangle, rounded corners=5pt,
      draw=gray!55, line width=0.95pt,
      fill=gray!5, minimum width=12.95cm, minimum height=1.22cm
    },
    note/.style={
      rectangle, rounded corners=4pt,
      draw=gray!45, line width=0.9pt,
      fill=yellow!9, minimum width=9.9cm, minimum height=0.66cm
    }
  ]

    % Column titles
    \node[font=\bfseries\small, text=blue!68!black]   at (-4.95,3.06) {FP16 注意力流程};
    \node[font=\bfseries\small, text=red!72!black]    at (-0.42,3.06) {量化模块};
    \node[font=\bfseries\small, text=orange!82!black] at (4.55,3.06) {耦合传播路径};

    % Left column
    \node[fp] (in)     at (-4.95,2.25) {\((q, K, V)\)};
    \node[fp] (logits) at (-4.95,1.20) {\(z = qK^\top / \sqrt{d_k}\)};
    \node[fp] (attn)   at (-4.95,0.15) {\(a = \mathrm{softmax}(z)\)};
    \node[fp] (out)    at (-4.95,-0.90) {\(o = aV\)};

    \draw[arr, blue!68!black] (in.south) -- (logits.north);
    \draw[arr, blue!68!black] (logits.south) -- (attn.north);
    \draw[arr, blue!68!black] (attn.south) -- (out.north);

    % Middle column
    \node[font=\bfseries, text=red!72!black] at (-0.42,2.37) {\(Q(\cdot)\)};
    \node[quant] (kq) at (-0.42,1.45) {\(K \rightarrow \hat K\)};
    \node[quant] (vq) at (-0.42,-0.10) {\(V \rightarrow \hat V\)};

    \draw[qarr] ([xshift=2pt,yshift=5pt]in.east)  to[bend left=11]  (kq.west);
    \draw[qarr] ([xshift=2pt,yshift=-5pt]in.east) to[bend right=11] (vq.west);

    % Right column
    \node[kpath] (kp) at (4.55,1.55) {
      \textbf{Key 误差 \(\rightarrow\) logits 偏移}\\[-1pt]
      \(\rightarrow\) 分布扭曲 \(\rightarrow \hat a-a\)
    };
    \node[vpath] (vp) at (4.55,-0.20) {
      \textbf{Value 误差 \(\rightarrow\) 聚合内容偏移}\\[-1pt]
      \(\rightarrow\) 聚合侧传播项
    };

    \draw[arr, cyan!60!black]   (kq.east) -- (kp.west);
    \draw[arr, orange!82!black] (vq.east) -- (vp.west);

    % Bottom decomposition
    \node[eqcard] (eq) at (0.15,-2.12) {};
    \node[font=\footnotesize] at (0.15,-2.12) {%
      \textbf{输出误差分解：}\quad
      $\textstyle
      \hat o-o
      =
      \big[\sum_{i=1}^{S}\hat a_i(\hat v_i-v_i)\big]_{\text{Value 聚合项}}
      +
      \big[\sum_{i=1}^{S}(\hat a_i-a_i)v_i\big]_{\text{Key 分布项}}
      $%
    };

    % subtle guide marks
    \draw[blue!18, line width=0.6pt] (-3.15,2.75) -- (-3.15,-1.25);
    \draw[red!18, line width=0.6pt] (1.05,2.75) -- (1.05,-1.25);

  \end{tikzpicture}
  \caption{注意力误差的两条传播路径及其代数分解。FP16 attention 给出参考分布 \(a\) 与输出 \(o\)，量化后的 Key 误差经 logits 与 softmax 改变注意力分布，Value 误差在量化后权重下进入加权聚合。底部公式给出式~\eqref{eq:ch3-error-decomp} 的最终分解形式。}
  \label{fig:error-decomposition}
\end{figure}


上述分解说明，下游任务感知到的是注意力行为偏移，而不是 $K$ 或 $V$ 的单独数值距离。Key 误差先改变 logits 的相对排序，再经 softmax 重分配概率质量；Value 误差改变被聚合的信息内容，并由量化后的分布 $\hat a$ 加权。只看张量重建误差，容易漏掉 softmax 非线性和加权聚合共同造成的功能失真。

本节由单头 Decode 注意力出发，将 KV Cache 量化形式化为注意力行为保持问题。第 \ref{sec:ch3-motivation-kv} 节据此诊断低比特场景下哪一侧更先触发失稳，第 \ref{sec:ch3-framework}--\ref{sec:ch3-allocator} 节再把这一目标落到离线校准产物、路径实例和预算分配机制中。

\section{动机诊断：K/V 敏感性不对称及其设计启示}
\label{sec:ch3-motivation-kv}

第 \ref{sec:ch3-problem} 节表明，Key 与 Value 的量化误差会分别沿注意力分布扭曲与加权聚合偏移两条路径共同作用于输出，但二者在低比特场景下未必同样脆弱。对于后续格式设计，关键问题更具体：在受限位宽预算下，究竟是哪一侧更先触发注意力行为的系统性失稳。只有先回答这一问题，后续低比特路径的格式选择才具有明确的结构性依据，而不会退化为针对具体模型或具体任务的经验调参。

为突出低比特失稳的角色依赖性，本文从第 \ref{sec:exp-kv-sensitivity} 节的完整证据中抽取图~\ref{fig:ch3-kv-asymmetry}。图中 \texttt{FP16} 是未量化参考，\texttt{K8V8} 是高精度量化参考，\texttt{K8V4} 与 \texttt{K4V8} 构成同平均 bit 预算下的角色对照，\texttt{K4V4} 则提供对称低比特锚点。其中，\texttt{K8V4} 与第四章完整诊断中的恢复性配置 \texttt{MixedKV} 对应。这个压缩视图只回答一个方法问题：在受限位宽预算下，4-bit 应优先落在 Key 还是 Value；完整跨模型读数和统计口径仍以后文表~\ref{tab:ch4-kv-ppl}、表~\ref{tab:ch4-kv-multitask} 与图~\ref{fig:ch4-kv-ruler32} 为准。

这组诊断把问题从“量化会不会带来误差”收窄到“4-bit 落在何侧”。在 Qwen2.5-1.5B 的单侧 PPL 诊断中，仅压低 Key 精度的 \texttt{K@INT4+V@FP16} 将 PPL 从 9.31 推到 1290.9，而仅压低 Value 精度的 \texttt{K@FP16+V@INT4} 为 9.35，见表~\ref{tab:ch4-kv-ppl}。在 32K 任务诊断中，\texttt{K4V8} 使 Qwen2.5-1.5B 与 Qwen2.5-7B 的 RULER 通过率降至 0\%；\texttt{MixedKV}（图中 \texttt{K8V4}）未出现这种归零式失稳，但其恢复强度仍随模型和任务变化，见表~\ref{tab:ch4-kv-multitask}。这些锚点说明，对当前低比特设定而言，Qwen 系列中首先越过功能临界点的是 Key 精度下降，而不是 Value 的同步压缩。

% 图 3-2 建议置于此处
% 图题：图 3-2 K/V 量化敏感性不对称的压缩版诊断与设计启示
% 数据纪律：本图必须是第 4.3.2 节同源完整证据中抽取出的位宽布局压缩视图；右面板只保留同口径 Needle，RULER 留在第 4.3.2 节完整证据中
\begin{figure}[!htbp]
  \centering
  \resizebox{\linewidth}{!}{%
  \begin{tikzpicture}[
    x=1cm, y=1cm,
    every node/.style={font=\scriptsize, align=center},
    panel/.style={
      rectangle, rounded corners=5pt,
      draw=black!25, fill=black!1, line width=0.7pt
    },
    title/.style={font=\bfseries\small, anchor=west, text=black},
    note/.style={font=\scriptsize, anchor=west, text=black!64},
    statebox/.style={
      rectangle, rounded corners=3pt,
      minimum width=0.72cm, minimum height=0.43cm,
      inner sep=2.1pt, line width=0.75pt
    },
    chainstep/.style={
      rectangle, rounded corners=3pt,
      text width=4.55cm, minimum height=0.48cm,
      inner sep=3.6pt, draw=black!38, fill=white, line width=0.65pt
    },
    arr/.style={-stealth, draw=black!58, line width=0.75pt}
  ]
    \definecolor{kBlue}{RGB}{58,104,173}
    \definecolor{kBlueFill}{RGB}{224,235,249}
    \definecolor{vTeal}{RGB}{66,139,133}
    \definecolor{vTealFill}{RGB}{224,242,239}
    \definecolor{riskRed}{RGB}{184,84,78}
    \definecolor{riskRedFill}{RGB}{249,228,225}
    \definecolor{anchorGold}{RGB}{178,130,26}
    \definecolor{anchorGoldFill}{RGB}{250,239,202}
    \definecolor{safeGreen}{RGB}{73,135,96}
    \definecolor{safeGreenFill}{RGB}{230,243,234}

    % panel backgrounds
    \node[panel, minimum width=8.65cm, minimum height=5.65cm, anchor=south west] at (0.00,0.00) {};
    \node[panel, minimum width=5.85cm, minimum height=5.65cm, anchor=south west] at (8.95,0.00) {};

    % panel (a)
    \node[title] at (0.30,5.30) {（a）位宽落点诊断};
    \node[note] at (0.30,4.95) {同平均位宽预算下，4bit 应优先落在 K 还是 V？};

    \node[font=\tiny\bfseries, text=black!70] at (1.85,4.42) {8 位基准};
    \draw[black!28, line width=0.55pt] (1.28,4.26) -- (2.42,4.26);
    \node[font=\tiny\bfseries, text=black!70] at (3.95,4.42) {同预算对照};
    \draw[black!28, line width=0.55pt] (2.75,4.26) -- (5.15,4.26);
    \node[font=\tiny\bfseries, text=black!70] at (6.25,4.42) {对称低比特锚点};
    \draw[black!28, line width=0.55pt] (5.60,4.26) -- (6.90,4.26);

    \node[font=\scriptsize\bfseries] at (1.85,3.94) {\texttt{K8V8}};
    \node[font=\scriptsize\bfseries] at (3.25,3.94) {\texttt{K8V4}};
    \node[font=\scriptsize\bfseries] at (4.65,3.94) {\texttt{K4V8}};
    \node[font=\scriptsize\bfseries] at (6.25,3.94) {\texttt{K4V4}};

    \node[anchor=east, font=\tiny\bfseries, align=right] at (1.05,3.22) {Qwen\\1.5B};
    \node[anchor=east, font=\tiny\bfseries, align=right] at (1.05,2.42) {Qwen\\7B};
    \node[anchor=east, font=\tiny\bfseries, align=right] at (1.05,1.62) {LLaMA\\8B};

    \foreach \y in {2.82,2.02,1.22} {
      \draw[black!10] (1.28,\y) -- (7.10,\y);
    }

    \node[statebox, fill=safeGreenFill, draw=safeGreen, text=safeGreen!65!black] at (1.85,3.22) {可用};
    \node[statebox, fill=safeGreenFill, draw=safeGreen, text=safeGreen!65!black] at (3.25,3.22) {可用};
    \node[statebox, fill=riskRedFill, draw=riskRed, text=riskRed!80!black] at (4.65,3.22) {崩塌};
    \node[statebox, fill=riskRedFill, draw=riskRed, text=riskRed!80!black] at (6.25,3.22) {崩塌};

    \node[statebox, fill=safeGreenFill, draw=safeGreen, text=safeGreen!65!black] at (1.85,2.42) {可用};
    \node[statebox, fill=safeGreenFill, draw=safeGreen, text=safeGreen!65!black] at (3.25,2.42) {可用};
    \node[statebox, fill=riskRedFill, draw=riskRed, text=riskRed!80!black] at (4.65,2.42) {崩塌};
    \node[statebox, fill=riskRedFill, draw=riskRed, text=riskRed!80!black] at (6.25,2.42) {崩塌};

    \node[statebox, fill=safeGreenFill, draw=safeGreen, text=safeGreen!65!black] at (1.85,1.62) {可用};
    \node[statebox, fill=safeGreenFill, draw=safeGreen, text=safeGreen!65!black] at (3.25,1.62) {可用};
    \node[statebox, fill=safeGreenFill, draw=safeGreen, text=safeGreen!65!black] at (4.65,1.62) {保留};
    \node[statebox, fill=safeGreenFill, draw=safeGreen, text=safeGreen!65!black] at (6.25,1.62) {保留};

    \node[font=\tiny, text=black!62, align=left, anchor=west] at (0.65,0.48)
      {代表性检索观察（Needle）：K8V8 与 K8V4 保持可用；\\Qwen 系列在 4bit K 下更易崩塌；LLaMA-3.1-8B 保留较高通过率，\\说明强度受架构调制。};

    % panel (b)
    \node[title] at (9.25,5.30) {（b）设计启示链};
    \node[note] at (9.25,4.95) {从诊断观察收束为低比特路径约束};

    \node[chainstep, fill=kBlueFill, draw=kBlue] (c1) at (11.88,4.28)
      {高精度路径保持可用};
    \node[chainstep, fill=vTealFill, draw=vTeal] (c2) at (11.88,3.34)
      {同预算落点暴露失效侧};
    \node[chainstep, fill=riskRedFill, draw=riskRed] (c3) at (11.88,2.40)
      {4bit K 可能越过注意力排序临界点};
    \node[chainstep, fill=anchorGoldFill, draw=anchorGold] (c4) at (11.88,1.42)
      {优先保护 K，再区分 K/V 角色};
    \node[chainstep, fill=black!4, draw=black!45] (c5) at (11.88,0.52)
      {角色感知非对称格式};

    \draw[arr] (c1.south) -- (c2.north);
    \draw[arr] (c2.south) -- (c3.north);
    \draw[arr] (c3.south) -- (c4.north);
    \draw[arr] (c4.south) -- (c5.north);
  \end{tikzpicture}%
  }
  \caption{K/V 位宽落点诊断的压缩版方法动机。子图（a）将 \texttt{K8V8} 基准、同平均位宽预算下的 \texttt{K8V4}/\texttt{K4V8} 对照，以及 \texttt{K4V4} 对称低比特锚点合并为 $3\times4$ 诊断矩阵，只保留代表性 Needle 观察：\texttt{K8V8} 与 \texttt{K8V4} 保持可用，Qwen 系列在 \texttt{K4V8}/\texttt{K4V4} 下出现 K 侧崩塌，而 LLaMA-3.1-8B 不呈现同强度坍塌。子图（b）将该观察收束为设计启示链：高精度路径保持可用，同预算落点暴露失效侧，4bit K 可能越过注意力排序临界点，因此低比特路径应优先保护 K，并进一步区分 K/V 角色。本图不承担完整实验比较；完整 Needle、RULER 与 LongBench 风格证据见第~\ref{subsec:ch4-kv-diagnosis}~节。}
  \label{fig:ch3-kv-asymmetry}
\end{figure}


进一步地，这种不对称性在 Qwen 系列上表现得最尖锐，在 LLaMA-3.1-8B 上则受到模型家族、规模和 GQA 配置的调制。本文的代表性诊断表明，Qwen 系列在 \texttt{K4V8} 与 \texttt{K4V4} 下更容易出现近乎完全的检索失稳，而具有更高 $H_{kv}$ 或不同头共享结构的模型则表现出更强的鲁棒性。本文据此把“Key 主导退化”限定为低比特路径设计中的结构性风险提示：方法需要优先避免 Key 侧过早跌破功能临界点，同时允许不同模型家族、模型规模与 GQA 结构在恢复强度上呈现差异。

基于上述动机诊断，本文将低比特路径的设计目标从“平均压缩 K 与 V”转化为“避免让 4-bit 率先落在 Key 上，并优先保护 Key 的表示能力”。在 Qwen 系列上，如果 \texttt{K4V8} 与 \texttt{K4V4} 的共同失效都指向 Key 精度下降，继续沿用对 K 与 V 一视同仁的对称压缩便缺乏结构依据；更合理的方向，是让量化格式本身对二者的不同功能角色作出响应。基于这一判断，第 \ref{sec:ch3-framework} 节将上述判断转成模块接口，第 \ref{sec:ch3-paths} 节则把这一设计动机实例化为从对称 INT4 向角色感知低比特路径的格式升级，并最终导向 \texttt{INT4-RoleAlign} 的正式定义。

\section{行为引导量化框架总览}
\label{sec:ch3-framework}
\label{sec:ch3-overview}

前两节给出两个约束：量化质量应按注意力行为读，而低比特失稳又对 Key 与 Value 不对称。本节把这两个约束转成后续模块之间的接口。离线校准模块接收校准样本、候选量化路径和可行域约束，输出路径特定的量化产物；预算分配模块读取同一离线链路中的逐层行为读数，输出逐层或逐角色的 bit 决策。这样写的目的，是先说明每个模块消费什么、产出什么，再讨论具体路径如何实例化。

形式上，本文将逐层行为敏感度画像记为
\begin{equation}
\mathcal{S}=\{s^{(l)}\}_{l=1}^{L},
\end{equation}
其中 $s^{(l)}$ 表示第 $l$ 层对量化扰动的行为敏感度读数。$\mathcal{S}$ 不预设唯一计算公式；它只规定一个接口语义：某层的读数越高，该层在预算受限时越应被谨慎压缩。校准层的行为偏移由路径相关代理计算，分配层则从同一离线校准链路中读取可聚合统计量，形成逐层敏感度画像。二者共享离线来源和行为读法，不要求数值表达式逐项相同。

在框架的第一层，即校准层，问题被组织为参数选择问题。设候选量化参数空间为 $\Theta$，则校准层的目标可抽象写为
\begin{equation}
\theta^\star = \arg\min_{\theta\in\Theta}\Delta_{\mathrm{beh}}^{\mathrm{cal}}(\theta),
\end{equation}
其中 $\Delta_{\mathrm{beh}}^{\mathrm{cal}}$ 是校准层实际使用的行为代理。给定某条路径及其候选参数家族，校准层输出 $\theta^\star$，供在线缓存写入模块直接读取；不同位宽和格式的差异体现在候选参数家族与代理定义上。第 \ref{sec:ch3-calibration} 节先定义这一离线选择规则，第 \ref{sec:ch3-paths} 节再分别落到 \texttt{INT8} 基准路径与低比特角色感知路径。

在框架的第二层，即分配层，问题被组织为预算决策问题。设平均 KV bit 预算为 $\bar b$，则分配器可抽象写为
\begin{equation}
b^\star = \mathcal{A}(\mathcal{S};\bar b),
\end{equation}
其中 $b^\star$ 是逐层或逐角色预算向量。分配层不重新搜索量化参数，而是读取 $\mathcal{S}$ 并回答“哪里保留更高精度”。最基础的输出是逐层预算；进一步的输出可以包含 K/V 角色条件化预算和 \texttt{AutoK} 建议。本文把 allocation / \texttt{AutoK} 写成校准产物的下游消费者，而不是与校准层并列的另一套目标函数。

执行时只有两个阶段。离线阶段提取参考行为、搜索路径参数、写出校准产物，并生成逐层敏感度读数；在线阶段读取这些固定产物，并按已冻结规则计算当前缓存写入所需的尺度或预算，不重新组织目标函数，也不在推理过程中重复离线搜索。这个分工保证后续系统实现面对的是冻结参数、规则和元数据，而不是一段需要在线重跑的分析过程。

% 图 3-3 建议置于此处
% 图题：图 3-3 行为引导量化框架总览：校准与分配共享同一行为敏感度画像
\begin{figure}[!htbp]
  \centering
  \resizebox{0.98\linewidth}{!}{%
  \begin{tikzpicture}[
    x=1cm, y=1cm,
    >=Stealth,
    every node/.style={font=\rmfamily\small, align=center, text=black},
    stage/.style={
      rectangle, rounded corners=6pt,
      draw=black!22, fill=black!1, line width=0.65pt
    },
    stageLabel/.style={font=\bfseries\scriptsize, text=black!68, anchor=west},
    card/.style={
      rectangle, rounded corners=5pt,
      draw=black!28, fill=white, line width=0.72pt,
      minimum height=3.15cm
    },
    smallcard/.style={
      rectangle, rounded corners=5pt,
      draw=black!28, fill=white, line width=0.72pt,
      minimum height=1.28cm
    },
    cardTitle/.style={font=\bfseries\scriptsize, text=black!78},
    row/.style={
      rectangle, rounded corners=3pt,
      minimum height=0.46cm, inner sep=2.3pt,
      font=\scriptsize, line width=0.55pt
    },
    note/.style={font=\scriptsize, text=black!58},
    arr/.style={-Stealth, draw=black!50, line width=0.74pt},
    arrBlue/.style={-Stealth, draw=blueTone!72!black, line width=0.78pt},
    arrGreen/.style={-Stealth, draw=greenTone!70!black, line width=0.78pt},
    arrGold/.style={-Stealth, draw=goldTone!78!black, line width=0.78pt}
  ]
    \definecolor{blueTone}{RGB}{62,103,167}
    \definecolor{blueFill}{RGB}{230,238,250}
    \definecolor{greenTone}{RGB}{66,134,108}
    \definecolor{greenFill}{RGB}{229,243,236}
    \definecolor{goldTone}{RGB}{178,132,37}
    \definecolor{goldFill}{RGB}{251,240,211}
    \definecolor{grayFill}{RGB}{245,246,248}

    \node[stage, minimum width=13.80cm, minimum height=5.42cm] (offline) at (6.95,2.76) {};
    \node[stage, minimum width=3.18cm, minimum height=5.42cm] (onlineStage) at (15.70,2.76) {};
    \node[stageLabel] at (0.34,5.28) {离线阶段};
    \node[stageLabel] at (14.30,5.28) {在线推理阶段};

    \node[card, minimum width=2.52cm] (input) at (1.55,2.72) {};
    \node[cardTitle] at (1.55,4.03) {离线输入};
    \node[row, draw=black!24, fill=grayFill, minimum width=1.82cm] at (1.55,3.45) {校准样本};
    \node[row, draw=black!24, fill=grayFill, minimum width=1.82cm] at (1.55,2.72) {FP16 参考行为};
    \node[row, draw=black!24, fill=grayFill, minimum width=1.82cm] at (1.55,1.99) {候选路径 $\Theta$};
    \node[note] at (1.55,1.31) {固定校准集};

    \node[card, draw=blueTone!60!black, fill=blueFill, minimum width=3.18cm] (calib) at (4.95,2.72) {};
    \node[cardTitle, text=blueTone!64!black] at (4.95,4.03) {离线校准链路};
    \node[row, draw=blueTone!55!black, fill=white, minimum width=2.34cm] at (4.95,3.43)
      {计算 $\Delta_{\mathrm{beh}}^{\mathrm{cal}}$};
    \node[row, draw=blueTone!55!black, fill=white, minimum width=2.34cm] at (4.95,2.72)
      {候选搜索 $\theta\in\Theta$};
    \node[row, draw=blueTone!55!black, fill=white, minimum width=2.34cm] at (4.95,2.01)
      {稳健选择};
    \node[note] at (4.95,1.31) {路径参数在离线侧确定};

    \node[smallcard, draw=blueTone!60!black, fill=white, minimum width=2.72cm] (theta) at (8.75,3.72) {};
    \node[cardTitle, text=blueTone!64!black] at (8.75,4.08) {冻结校准产物};
    \node[font=\bfseries\small] at (8.75,3.68) {$\theta^\star$};
    \node[note] at (8.75,3.30) {尺度、零点、角色格式};

    \node[smallcard, draw=greenTone!62!black, fill=greenFill, minimum width=2.72cm] (profile) at (8.75,1.72) {};
    \node[cardTitle, text=greenTone!58!black] at (8.75,2.09) {行为敏感度画像};
    \node[font=\bfseries\small] at (8.75,1.69) {$\mathcal S=\{s^{(l)}\}_{l=1}^{L}$};
    \node[note] at (8.75,1.30) {校准链路副产物};

    \node[smallcard, draw=goldTone!70!black, fill=goldFill, minimum width=2.82cm] (alloc) at (12.03,1.72) {};
    \node[cardTitle, text=goldTone!72!black] at (12.03,2.09) {预算分配};
    \node[font=\bfseries\small] at (12.03,1.69) {$\mathcal A(\mathcal S;\bar b)$};
    \node[note] at (12.03,1.30) {生成位宽向量 $b^\star$};

    \node[card, draw=black!32, fill=white, minimum width=2.72cm] (online) at (15.70,2.72) {};
    \node[cardTitle] at (15.70,4.03) {在线只读执行};
    \node[row, draw=blueTone!45!black, fill=blueFill, minimum width=1.82cm] at (15.70,3.45)
      {读取 $\theta^\star$};
    \node[row, draw=goldTone!55!black, fill=goldFill, minimum width=1.82cm] at (15.70,2.72)
      {读取 $b^\star$};
    \node[row, draw=black!24, fill=grayFill, minimum width=1.82cm] at (15.70,1.99)
      {缓存写入与解码};
    \node[note] at (15.70,1.31) {固定产物执行};

    \draw[arr] (input.east) -- node[above, font=\tiny, text=black!55] {参考与候选} (calib.west);

    \coordinate (fork) at ($(calib.east)+(0.42,0)$);
    \draw[arrBlue] (calib.east) -- (fork) |- node[pos=0.72, above, font=\tiny, text=blueTone!66!black] {冻结产物} (theta.west);
    \draw[arrGreen] (fork) |- node[pos=0.72, below, font=\tiny, text=greenTone!58!black] {副产物} (profile.west);
    \draw[arrGold] (profile.east) -- node[above, font=\tiny, text=goldTone!70!black] {只读画像} (alloc.west);

    \draw[arrBlue] (theta.east) -- ++(0.52,0) |- node[pos=0.70, above, font=\tiny, text=blueTone!66!black] {缓存参数} (online.west |- theta.east);
    \draw[arrGold] (alloc.east) -- ++(0.52,0) |- node[pos=0.70, below, font=\tiny, text=goldTone!70!black] {预算配置} (online.west |- alloc.east);
  \end{tikzpicture}%
  }
  \caption{行为引导量化框架总览。离线校准链路以校准样本、FP16 参考行为与候选量化路径为输入，计算行为代理并完成稳健选择，随后形成冻结校准产物 \(\theta^\star\)。同一链路同步生成逐层行为敏感度画像 \(\mathcal S\)，预算分配模块只读该画像，在平均位宽预算 \(\bar b\) 下生成位宽向量 \(b^\star\)。在线推理阶段读取 \(\theta^\star\) 与 \(b^\star\)，执行缓存写入与解码。}
  \label{fig:ch3-framework-shared-profile}
\end{figure}


后续各节按这个接口展开：第 \ref{sec:ch3-calibration} 节定义校准层代理与选择准则；第 \ref{sec:ch3-paths} 节给出 \texttt{INT8} 基准路径、对称 \texttt{INT4} 的局限和 \texttt{INT4-RoleAlign}；第 \ref{sec:ch3-allocator} 节把同源敏感度画像转成逐层预算与 \texttt{AutoK} 建议；第 \ref{sec:ch3-deployment} 节说明离线产物如何被在线推理管线读取。

\section{行为引导校准目标与参数搜索策略}
\label{sec:ch3-calibration}

第 \ref{sec:ch3-problem} 节给出了联合行为目标，但该目标不能直接作为离线搜索中的单一损失。本节把它转写为校准层可执行的代理：先定义注意力分布侧 KL，再说明它与 MSE、反向 KL 的取舍，最后给出参数搜索与选择纪律。

本节统一的是流程和产物边界，不是强行要求所有路径共享同一个评分函数。对称路径与角色感知低比特路径可以有不同候选参数家族和代理定义；共同点是都在离线样本上计算行为读数，在可行域约束下优先控制尾部坏案例，并把最终选择写入固定校准产物。

\subsection{注意力分布 KL 散度目标}

设第 $l$ 层、第 $h$ 个注意力头、在时间步 $t$ 的 Query、Key 分别记为 $q^{(l,h,t)}$ 与 $K^{(l,h,t)}$。在 FP16 参考路径下，其注意力分布写为
\begin{equation}
p_{\mathrm{ref}}^{(l,h,t)}
=
\operatorname{softmax}\!\left(
\frac{q^{(l,h,t)} K^{(l,h,t)\top}}{\sqrt{d_k}}
\right).
\end{equation}
对给定候选量化参数 $\theta$，记量化—反量化后的 Key 为 $\tilde K_{\theta}^{(l,h,t)}$，则对应的量化后注意力分布为
\begin{equation}
p_{\theta}^{(l,h,t)}
=
\operatorname{softmax}\!\left(
\frac{q^{(l,h,t)} \tilde K_{\theta}^{(l,h,t)\top}}{\sqrt{d_k}}
\right).
\end{equation}
基于此，本文将单个位置上的分布侧行为偏移代理定义为
\begin{equation}
d_{\mathrm{KL}}^{(l,h,t)}(\theta)
=
D_{\mathrm{KL}}
\!\left(
p_{\mathrm{ref}}^{(l,h,t)}
\;\middle\|\;
p_{\theta}^{(l,h,t)}
\right).
\end{equation}
本文优化的是\emph{分布侧 KL 代理}，不声称输出误差 $(\hat o-o)$ 与 KL 数学等价。KL 的作用是把联合行为目标落到校准层可执行的分布读数上。

为什么不直接用 MSE？原因在于 MSE 工作在张量重建空间，只比较 $K$ 或 $V$ 的逐元素差异；而 Key 误差真正进入模型功能前，要先经过 $qK^\top$、softmax 排序和概率质量分配。一个 MSE 较小的候选仍可能把关键 token 的注意力质量移走。选择 attention-distribution KL，是因为它直接比较参考分布与量化后分布。回顾
\[
\hat o-o
=
\sum_i(\hat a_i-a_i)v_i
+
\sum_i \hat a_i(\hat v_i-v_i),
\]
其中左项是分布侧传播路径，右项是聚合侧传播路径。KL 不能覆盖全部输出误差，但它直接约束 $(\hat a-a)$，也会影响聚合侧使用的权重 $\hat a_i$。因此，在 Key-sensitive 校准中，KL 是比逐元素重建误差更贴近注意力行为的代理。

为什么使用前向 KL？本文关心的是参考路径中高概率位置是否被保住。前向 KL $D_{\mathrm{KL}}(p_{\mathrm{ref}}\|p_\theta)$ 对参考高概率位置被量化路径低估的情况惩罚更强，适合长上下文检索中“不能漏掉关键 token”的需求。反向 KL 更偏向量化分布自身的高概率区域，JS 散度则更对称；它们可以作为补充诊断，但不是本文主路径的默认校准目标。

数值实现上，本文先对概率做小常数 $\varepsilon$ 截断，再计算 KL，以避免极端尾部概率造成不稳定。KL 在不同模型规模和 bit-width 下的收益强弱属于第四章的实验问题，不作为本节前提。另一个边界是：attention-distribution KL 服务于 Key-sensitive 的分布侧校准；Value 路径使用独立输出扰动代理，不被强行写成同一个 KL 目标。

\subsection{参数搜索空间与稳健选择准则}
\label{subsec:ch3-two-stage}

在上述代理定义之上，本文将离线校准组织为一个路径相关、但流程统一的参数搜索问题。给定某一路径对应的候选参数集合 $\Theta_{\mathrm{path}}$，其最优参数可抽象写为
\begin{equation}
\theta_{\mathrm{path}}^\star
=
\arg\min_{\theta \in \Theta_{\mathrm{path}}}
R_{\mathrm{path}}(\theta),
\end{equation}
其中 $R_{\mathrm{path}}(\theta)$ 表示该路径上的稳健风险统计量。这里的“统一”体现在工作流与纪律层面：都需要在校准样本上生成参考行为、计算路径对应的行为代理、聚合尾部与整体统计量、并在可行域约束下给出最终参数；但这并不要求所有路径都必须共享同一个标量 score function。对称路径可直接用 attention-KL 定义风险，而 role-aware low-bit path 则允许在 K 与 V 两侧使用不同的行为代理。

对称路径（包括 INT8 基准路径与对称 \texttt{INT4} 试探路径）的选择纪律可写为一组统一的稳健统计。设校准样本集合为 $\mathcal D_{\mathrm{calib}}$，则可定义整体均值、尾部统计与 K/V 两侧的裁剪率：
\begin{equation}
\mu(\theta)
=
\operatorname{Mean}_{(x,l,h,t)\in\mathcal D_{\mathrm{calib}}}
d_{\mathrm{KL}}^{(l,h,t)}(\theta),
\end{equation}
\begin{equation}
q_{0.95}(\theta)
=
\operatorname{P95}_{(x,l,h,t)\in\mathcal D_{\mathrm{calib}}}
d_{\mathrm{KL}}^{(l,h,t)}(\theta),
\end{equation}
\begin{equation}
c_K(\theta)=\mathrm{clip\mbox{-}rate}_K(\theta), \qquad
c_V(\theta)=\mathrm{clip\mbox{-}rate}_V(\theta).
\end{equation}
其中，$\mu(\theta)$ 反映整体平均行为偏移，$q_{0.95}(\theta)$ 反映尾部坏案例，而 $c_K(\theta)$ 与 $c_V(\theta)$ 用于约束“通过过度裁剪换取表面更低 KL”的伪优候选。基于此，首先定义满足双侧裁剪约束的可行域
\begin{equation}
\Theta_{\mathrm{feasible}}
=
\{\theta \in \Theta_{\mathrm{path}}: c_K(\theta)\le \tau_K,\; c_V(\theta)\le \tau_V\},
\end{equation}
再在可行域中按尾部稳健性优先选择
\begin{equation}
\theta^\star
=
\arg\min_{\theta \in \Theta_{\mathrm{feasible}}}
q_{0.95}(\theta).
\end{equation}
若所有候选均不满足裁剪率约束，则退化为优先最小化总裁剪负担 $c_K(\theta)+c_V(\theta)$，再以尾部统计与平均 KL 作为后续排序依据。该规则体现了本文一贯的校准纪律：不是简单选择“均值最小”的候选，而是优先压制长上下文场景中更具破坏性的尾部异常。

对于 \texttt{INT4-RoleAlign}，离线搜索不再被写成“单一 attention-KL 分数统一选出 $(p_K,p_V)$”，而是拆分为同一流程中的两条角色感知接口。其一，\textbf{K-path} 以 $p_K$ 为候选参数，在 per-channel Key 非对称量化格式上使用注意力分布 KL 作为主代理，并在同一套稳健搜索纪律下完成选择；其二，\textbf{V-path} 以 $p_V$ 为候选参数，在 per-token Value 非对称量化格式上采用独立的输出扰动代理。本文的正式方法定义将这一 V-path 代理与通用共享校准口径区分开来，避免把未区分 K/V 角色的校准目标直接等同为 \texttt{RoleAlign}。也就是说，对 \texttt{RoleAlign} 而言，变化的是各路径的参数家族与行为代理；不变的是离线校准流程、可行域约束、尾部稳健优先与产物输出接口。正因为如此，本文后续才能在同一框架下同时讨论对称路径与角色感知路径，而不必为了维持表面统一性而误把两条角色路径压成一个数学上并不存在的单一目标。

从输出形式上看，离线校准流程统一产生两类结果：一类是\emph{路径特定校准产物}，包括静态缩放参数、zero-point、分位数参数、分组元数据与保护参数等，用于在线推理时的直接读取与执行；不同路径对应的字段集合可以不同，但都服从同一只读产物接口。另一类是\emph{共享行为统计读数}，它们来自同一离线校准链路，并将在第 \ref{sec:ch3-allocator} 节中被进一步组织为逐层行为敏感度画像。校准层不仅回答“如何量化”，也为后续的预算分配提供同源读数。逐头温度校正等附加组件在本文主线中统一作为补充诊断对象，仅在附录中保留说明；主线方法默认不引入该组件。

% 图 3-4 建议置于本小节第 1 段与第 2 段之间
% 图题：图 3-4 行为引导离线校准的统一工作流
\begin{figure}[!htbp]
  \centering
  \resizebox{0.98\linewidth}{!}{%
  \begin{tikzpicture}[
    x=1cm, y=1cm,
    >=Stealth,
    every node/.style={font=\rmfamily\small, align=center, text=black},
    band/.style={
      rectangle, rounded corners=6pt,
      draw=black!18, fill=black!1, line width=0.62pt
    },
    card/.style={
      rectangle, rounded corners=5pt,
      draw=black!28, fill=white, line width=0.72pt,
      minimum height=3.18cm
    },
    cardTitle/.style={font=\bfseries\scriptsize, text=black!76},
    item/.style={
      rectangle, rounded corners=3pt,
      minimum height=0.44cm, inner sep=2.2pt,
      font=\scriptsize, line width=0.55pt
    },
    note/.style={font=\scriptsize, text=black!58},
    arr/.style={-Stealth, draw=black!48, line width=0.72pt},
    arrBlue/.style={-Stealth, draw=blueTone!72!black, line width=0.76pt},
    arrGold/.style={-Stealth, draw=goldTone!78!black, line width=0.76pt},
    arrGreen/.style={-Stealth, draw=greenTone!70!black, line width=0.76pt}
  ]
    \definecolor{tealTone}{RGB}{66,139,133}
    \definecolor{tealFill}{RGB}{228,243,240}
    \definecolor{blueTone}{RGB}{62,103,167}
    \definecolor{blueFill}{RGB}{230,238,250}
    \definecolor{goldTone}{RGB}{178,132,37}
    \definecolor{goldFill}{RGB}{251,240,211}
    \definecolor{greenTone}{RGB}{66,134,108}
    \definecolor{greenFill}{RGB}{229,243,236}
    \definecolor{grayFill}{RGB}{245,246,248}

    \node[band, minimum width=17.15cm, minimum height=4.72cm] at (8.78,2.62) {};
    \node[font=\bfseries\small, anchor=west] at (0.52,4.78) {行为引导离线校准工作流};
    \node[note, anchor=west] at (0.52,4.43) {每条量化路径共享同一选择纪律，差异落在候选网格、行为代理与产物字段};
    \draw[draw=black!14, line width=0.52pt] (0.52,4.20) -- (17.02,4.20);

    \node[card, draw=tealTone!70!black, fill=white, minimum width=2.58cm] (ref) at (1.70,2.20) {};
    \node[cardTitle, text=tealTone!64!black] at (1.70,3.46) {1 参考提取};
    \node[item, draw=tealTone!55!black, fill=tealFill, minimum width=1.80cm] at (1.70,2.86) {校准样本};
    \node[item, draw=tealTone!55!black, fill=white, minimum width=1.80cm] at (1.70,2.20) {FP16 前向};
    \node[item, draw=tealTone!55!black, fill=white, minimum width=1.80cm] at (1.70,1.54) {$p_{\mathrm{ref}},\,o_{\mathrm{ref}}$};
    \node[note] at (1.70,0.98) {参考行为};

    \node[card, draw=blueTone!62!black, fill=blueFill, minimum width=2.82cm] (grid) at (5.02,2.20) {};
    \node[cardTitle, text=blueTone!64!black] at (5.02,3.46) {2 路径实例化};
    \node[item, draw=blueTone!50!black, fill=white, minimum width=2.05cm] at (5.02,2.88) {$\Theta_{\mathrm{path}}$};
    \node[item, draw=blueTone!50!black, fill=white, minimum width=2.05cm] at (5.02,2.24) {INT8 / 对称 INT4};
    \node[item, draw=blueTone!50!black, fill=white, minimum width=2.05cm] at (5.02,1.60) {RoleAlign};
    \node[note] at (5.02,0.98) {候选网格};

    \node[card, draw=blueTone!62!black, fill=white, minimum width=2.88cm] (score) at (8.48,2.20) {};
    \node[cardTitle, text=blueTone!64!black] at (8.48,3.46) {3 行为评分};
    \node[item, draw=blueTone!50!black, fill=blueFill, minimum width=2.10cm] at (8.48,2.88) {$R_{\mathrm{path}}(\theta)$};
    \node[item, draw=blueTone!50!black, fill=white, minimum width=2.10cm] at (8.48,2.24) {K 侧 KL};
    \node[item, draw=blueTone!50!black, fill=white, minimum width=2.10cm] at (8.48,1.60) {V 侧输出扰动};
    \node[note] at (8.48,0.98) {路径相关代理};

    \node[card, draw=goldTone!70!black, fill=goldFill, minimum width=2.82cm] (select) at (11.94,2.20) {};
    \node[cardTitle, text=goldTone!75!black] at (11.94,3.46) {4 稳健选择};
    \node[item, draw=goldTone!60!black, fill=white, minimum width=2.05cm] at (11.94,2.88) {$\Theta_{\mathrm{feasible}}$};
    \node[item, draw=goldTone!60!black, fill=white, minimum width=2.05cm] at (11.94,2.24) {P95 尾部读数};
    \node[item, draw=goldTone!60!black, fill=white, minimum width=2.05cm] at (11.94,1.60) {裁剪率审计};
    \node[note] at (11.94,0.98) {选择纪律};

    \node[card, draw=greenTone!70!black, fill=greenFill, minimum width=2.72cm] (out) at (15.30,2.20) {};
    \node[cardTitle, text=greenTone!62!black] at (15.30,3.46) {5 产物冻结};
    \node[item, draw=greenTone!55!black, fill=white, minimum width=1.98cm] at (15.30,2.88) {$\theta_{\mathrm{path}}^\star$};
    \node[item, draw=greenTone!55!black, fill=white, minimum width=1.98cm] at (15.30,2.24) {行为读数 $\mathcal S$};
    \node[item, draw=greenTone!55!black, fill=white, minimum width=1.98cm] at (15.30,1.60) {在线接口};
    \node[note] at (15.30,0.98) {固定产物};

    \draw[arr] (ref.east) -- node[above, font=\tiny, text=black!55] {参考行为} (grid.west);
    \draw[arrBlue] (grid.east) -- node[above, font=\tiny, text=blueTone!62!black] {逐候选计算} (score.west);
    \draw[arrGold] (score.east) -- node[above, font=\tiny, text=goldTone!70!black] {风险统计} (select.west);
    \draw[arrGreen] (select.east) -- node[above, font=\tiny, text=greenTone!62!black] {冻结} (out.west);
  \end{tikzpicture}%
  }
  \caption{行为引导离线校准的统一工作流。校准样本先经过 FP16 前向得到参考行为；每条量化路径据此实例化候选网格 \(\Theta_{\mathrm{path}}\)，并计算路径相关风险 \(R_{\mathrm{path}}(\theta)\)。稳健选择阶段根据可行域、P95 尾部读数与裁剪率审计确定 \(\theta_{\mathrm{path}}^\star\)，同时保留行为读数 \(\mathcal S\) 与在线接口。RoleAlign 中，K 侧使用注意力分布 KL，V 侧使用输出扰动代理。}
  \label{fig:ch3-calibration-workflow}
\end{figure}


% 表 3-1 建议置于本小节第 3 段之后
% 表题：表 3-1 统一离线校准工作流下的路径参数化接口
\begin{table}[!htbp]
  \centering
  \caption{统一离线校准工作流下的路径参数化接口}
  \label{tab:ch3-calibration-interfaces}
  \scriptsize
  \setlength{\tabcolsep}{2.6pt}
  \renewcommand{\arraystretch}{1.12}
  \begin{tabularx}{\linewidth}{@{}>{\raggedright\arraybackslash}p{1.92cm}>{\raggedright\arraybackslash}p{1.38cm}>{\raggedright\arraybackslash}p{1.72cm}X>{\raggedright\arraybackslash}p{2.70cm}@{}}
    \toprule
    路径 & \makecell[l]{校准\\接口} & \makecell[l]{量化\\粒度} & \makecell[l]{主代理 / 约束} & \makecell[l]{输出\\产物} \\
    \midrule
    \makecell[l]{INT8 基准\\路径} & $(p_c,g,m)$ & \makecell[l]{对称\\per-group} & \makecell[l]{注意力分布 KL；\\P95；裁剪率约束} & \makecell[l]{静态 scale +\\保护元数据} \\
    \makecell[l]{对称 INT4\\试探路径} & $(p_c,g)$ & \makecell[l]{对称\\per-group} & \makecell[l]{注意力分布 KL；\\P95；裁剪率约束} & 静态 scale \\
    \makecell[l]{INT4-\\RoleAlign} & \makecell[l]{K: $p_K$\\V: $p_V$} & \makecell[l]{K: per-channel\\V: per-token} & \makecell[l]{K: 注意力分布 KL / P95\\V: 输出扰动代理\\路径相关稳健约束} & \makecell[l]{角色感知产物\\K/V 参数与轴元数据} \\
    \bottomrule
  \end{tabularx}

  \vspace{4pt}
  \parbox{0.96\textwidth}{\footnotesize 统一的是离线校准流程与稳健选择纪律，不同路径的差异体现在参数接口与行为代理上。在 \texttt{INT4-RoleAlign} 中，Key 路径由 attention-distribution KL 主导，Value 路径由独立的输出扰动代理主导。}
\end{table}


\section{跨位宽量化路径的实例化设计}
\label{sec:ch3-paths}

第 \ref{sec:ch3-calibration} 节已经将行为引导原则具体化为统一的离线校准流程与稳健选择纪律，但这一流程本身并不规定唯一的位宽或唯一的量化格式。对本文而言，真正需要回答的问题是：当位宽与格式约束发生变化时，同一行为原则如何被实例化为可执行的量化路径。基于这一思路，本节按照由保守到激进的顺序组织三条路径：首先给出 \texttt{INT8} 基准路径，作为行为引导校准最干净、最可审计的验证实例；随后说明直接将同一思路机械降到对称 \texttt{INT4} 为什么会遇到结构性困难；最后引入 \texttt{INT4-RoleAlign} 作为角色感知的低比特实例，并在方法层面将其与 \texttt{KIVI-style} 的设计边界明确区分开来。

需要强调的是，本节的目标不是在方法章内提前报告完整实验结果，而是说明同一框架如何在不同位宽与不同量化格式约束下落成具体路径。后续路径之间变化的是参数接口、量化轴与行为代理的具体定义，不变的是第 \ref{sec:ch3-calibration} 节已经建立的离线校准流程、可行域约束与稳健选择纪律。正是在这一点上，本文的方法实例既保持了跨位宽的一致性，也为后续 allocator 与系统落地提供了统一接口。

% 表 3-2 建议置于此处
% 表题：表 3-2 跨位宽路径的实例化关系与统一校准工作流
\begin{table}[!htbp]
  \centering
  \caption{跨位宽路径的实例化关系与统一校准工作流}
  \label{tab:ch3-path-instantiation}
  \small
  \setlength{\tabcolsep}{4.2pt}
  \renewcommand{\arraystretch}{1.12}
  \begin{tabular}{@{}>{\raggedright\arraybackslash}p{1.85cm}>{\raggedright\arraybackslash}p{2.05cm}>{\raggedright\arraybackslash}p{1.50cm}>{\raggedright\arraybackslash}p{4.05cm}>{\raggedright\arraybackslash}p{3.00cm}@{}}
    \toprule
    路径 & \makecell[l]{位宽 /\\量化格式} & \makecell[l]{主要\\接口} & \makecell[l]{暴露问题 / 目标} & \makecell[l]{统一 workflow\\中的定位} \\
    \midrule
    \makecell[l]{INT8 规范\\路径} & \makecell[l]{8-bit\\对称 per-group} & $(p_c,g,m)$ & 形成最干净、最可审计的规范验证实例 & \makecell[l]{规范验证路径；\\验证校准闭环} \\
    \makecell[l]{对称 INT4\\试探路径} & \makecell[l]{4-bit\\对称 per-group} & $(p_c,g)$ & \makecell[l]{暴露 Key 排序脆弱性，以及\\量化级别骤降带来的结构困难} & \makecell[l]{说明对称 INT4\\不能机械继承 INT8 路径} \\
    \makecell[l]{INT4-\\RoleAlign} & \makecell[l]{4-bit\\K: per-channel\\V: per-token} & \makecell[l]{K-path: $p_K$\\V-path: $p_V$} & \makecell[l]{在低比特下恢复 behavior 保持，\\并将 K/V 两侧纳入可审计接口} & \makecell[l]{角色感知低比特路径；\\统一 workflow 下的实例} \\
    \bottomrule
  \end{tabular}

  \vspace{4pt}
  \parbox{0.96\textwidth}{\footnotesize 注：统一的是离线校准 workflow 与稳健选择纪律；不同路径的差异体现在参数接口、量化格式与行为代理上。对 \texttt{INT4-RoleAlign} 而言，Key 路径由 attention-distribution KL 主导，Value 路径由独立的输出扰动代理主导。}
\end{table}


\subsection{INT8 基准路径：静态 Scale 与自适应保护}

本文首先在 \texttt{INT8} 位宽下给出行为引导校准的 INT8 基准路径，其目的不是追求极限压缩，而是建立一个最干净、最稳定、最便于审计的验证实例。相较于 \texttt{INT4}，\texttt{INT8} 提供更细的量化网格，因此更适合作为校准层闭环是否成立的第一性验证对象。也正因为如此，本文将该路径定位为 INT8 基准路径：它在方法论上承担“基准验证实例”的角色，而不是在压缩率层面承担“最终目标路径”的角色。

在该路径中，Key 与 Value 均采用对称、逐组（per-group）的静态量化。记第 $l$ 层第 $j$ 个分组内的待量化张量为 $x_{l,j}$，对给定校准百分位参数 $p_c$ 与分组大小 $g$，其静态缩放参数定义为
\begin{equation}
s_{l,j}^{\mathrm{static}}
=
\frac{\operatorname{Percentile}\!\big(|x_{l,j}|; p_c\big)}
{q_{\max}},
\end{equation}
其中 $q_{\max}$ 表示对称量化网格的正端点。由此，对应的量化—反量化写为
\begin{equation}
q
=
\operatorname{clamp}\!\left(
\operatorname{round}\!\left(\frac{x}{s_{l,j}^{\mathrm{static}}}\right),
-q_{\max},\, q_{\max}
\right),
\qquad
\hat x = q\, s_{l,j}^{\mathrm{static}}.
\end{equation}
这里真正由第 \ref{sec:ch3-calibration} 节确定的是 $(p_c,g)$ 如何被搜索与筛选；本小节关心的则是这些离线选出的参数如何被具体地转写为一条对称 \texttt{INT8} 路径。

在静态缩放参数之外，本文还为 INT8 基准路径引入自适应保护机制，用以处理运行时输入分布超出校准覆盖范围的情况。记由当前写入 token 估计得到的动态尺度为 $s_{l,j}^{\mathrm{dyn}}$，保护裕度参数为 $m$，则最终使用的缩放参数可写为
\begin{equation}
s_{l,j}^{\mathrm{final}}
=
\max\!\left(
m \cdot s_{l,j}^{\mathrm{static}},
\;
s_{l,j}^{\mathrm{dyn}}
\right).
\end{equation}
这一设计的作用不是把路径重新变回运行时校准，而是在离线校准主导的前提下，对极端分布漂移进行保守防护。静态缩放参数仍是主路径，自适应保护只是一个在必要时触发的安全阀，用于避免异常输入导致的剧烈裁剪。关于该机制在缓存写入时的具体系统语义——例如其是否会回写历史缓存——统一留到第 \ref{sec:ch3-deployment} 节再说明。

因此，INT8 基准路径在本文中的地位应被理解为：它是行为引导校准的基准实例，而不是极限压缩方案。它证明在最保守的位宽设置下，注意力行为可以被稳定地对齐和保持，并为后续更激进的低比特路径提供一个方法上可审计、实现上可复现的参照基线。

\subsection{对称 INT4 的局限与格式升级动因}

如果直接沿用上述对称路径，仅将位宽从 \texttt{INT8} 机械降至对称 \texttt{INT4}，则量化问题的性质将发生变化。这里的关键并不只是“误差更大了一些”，而是量化网格的分辨率发生了阶跃式下降：对称 \texttt{INT8} 具有更密的有效量化等级，而对称 \texttt{INT4} 仅保留极少数离散表示，这意味着原先在 \texttt{INT8} 下可由细粒度缩放参数调整吸收的分布差异，在 \texttt{INT4} 下会被更粗糙地投影到统一网格上。

这种网格粗化对 Key 的影响尤为敏感。第 \ref{sec:ch3-problem} 节已经说明，Key 扰动首先作用于 logits 的相对排序，并通过 softmax 放大或重排注意力分布；第 \ref{sec:ch3-motivation-kv} 节则进一步给出动机诊断：在当前出现 cliff 的 Qwen 设置中，首先触发对称 \texttt{INT4} 崩塌的证据指向 Key 精度下降，而非 Value 的同步退化。因此，当位宽降至 \texttt{INT4} 后，真正暴露出来的问题不再只是“数值重建精度降低”，而是对称量化格式本身无法充分适应 Key 通道的异质性，进而更容易破坏注意力排序结构。

这也意味着，后续设计不应再停留在“继续搜索更好的对称参数”这一层面，而必须把设计维度从参数优化提升到格式升级。如果导致对称 \texttt{INT4} 失效的主导因素来自 Key 侧的结构性脆弱，那么更合理的回应就不是对 $K$ 与 $V$ 继续施加相同的量化轴与相同的表示方式，而是让量化格式本身对二者在注意力计算中的不同功能角色作出响应。基于这一判断，下一小节将把低比特路径实例化为 \texttt{INT4-RoleAlign}：在保持第 \ref{sec:ch3-calibration} 节统一工作流不变的前提下，对 Key 与 Value 分别采用更契合其角色的非对称量化轴。

\subsection{INT4-RoleAlign：角色感知非对称量化}
\label{subsec:ch3-rolealign}
\label{sec:ch3-rolealign}

\texttt{INT4-RoleAlign} 是行为引导框架在低比特场景下的角色感知实例。其核心思想是：在保持统一离线校准流程与稳健选择纪律不变的前提下，不再要求 $K$ 与 $V$ 共享相同的量化轴，而是根据二者在注意力计算中的不同角色，分别采用不同的非对称量化形式。更具体地说，Key 侧优先保持注意力排序与通道异质性，Value 侧则在保持可聚合性的前提下适配时序动态范围变化。由此，\texttt{RoleAlign} 的关键不在于“比对称 \texttt{INT4} 多了一个公式”，而在于将第 \ref{sec:ch3-motivation-kv} 节识别出的角色不对称，转写为低比特路径中的量化轴不对称。

对 Key 而言，本文采用逐通道（per-channel）的非对称量化。设第 $l$ 层 Key 矩阵为 $K^{(l)}\in\mathbb{R}^{S\times d_k}$，则对第 $j$ 个通道，其分位数裁剪边界定义为
\begin{equation}
m_{l,j}^{K}
=
\operatorname{Percentile}\!\big(K^{(l)}_{:,j};\,100-p_K\big),
\qquad
M_{l,j}^{K}
=
\operatorname{Percentile}\!\big(K^{(l)}_{:,j};\,p_K\big),
\end{equation}
由此得到 Key 侧的缩放参数与仿射偏移项：
\begin{equation}
s_{l,j}^{K}
=
\frac{M_{l,j}^{K}-m_{l,j}^{K}}{q_{\max}-q_{\min}},
\qquad
o_{l,j}^{K}
=
m_{l,j}^{K}-q_{\min}s_{l,j}^{K}.
\end{equation}
这里的关键在于：分位数统计沿序列维度聚合、并对每个通道单独确定 $(s,o)$，从而更好地保持 Key 通道间的幅值异质性与相对排序结构。

对 Value 而言，本文采用逐 token（per-token）的非对称量化。设第 $l$ 层 Value 矩阵为 $V^{(l)}\in\mathbb{R}^{S\times d_v}$，则对第 $t$ 个 token，其分位数裁剪边界定义为
\begin{equation}
m_{l,t}^{V}
=
\operatorname{Percentile}\!\big(V^{(l)}_{t,:};\,100-p_V\big),
\qquad
M_{l,t}^{V}
=
\operatorname{Percentile}\!\big(V^{(l)}_{t,:};\,p_V\big),
\end{equation}
由此得到 Value 侧的缩放参数与仿射偏移项：
\begin{equation}
s_{l,t}^{V}
=
\frac{M_{l,t}^{V}-m_{l,t}^{V}}{q_{\max}-q_{\min}},
\qquad
o_{l,t}^{V}
=
m_{l,t}^{V}-q_{\min}s_{l,t}^{V}.
\end{equation}
此时，分位数统计沿特征维度聚合，并对每个 token 单独量化，从而使 Value 侧能够更灵活地适应不同时间步上的动态范围变化。

在 K/V 两侧，统一的量化—反量化形式均可写为
\begin{equation}
q
=
\operatorname{clamp}\!\left(
\operatorname{round}\!\left(\frac{x-o}{s}\right),
q_{\min},\,q_{\max}
\right),
\qquad
\hat x = q\,s+o.
\end{equation}
这里的 $o$ 表示仿射偏移项，用于和第二章标准非对称量化定义中的整数 zero-point 记号作区分。真正使 \texttt{RoleAlign} 成立的，并不是上述统一公式本身，而是参数 $(p_K,p_V)$ 的来源具有明确的角色划分：其中 $p_K$ 通过 per-channel Key 路径上的注意力分布 KL 搜索得到，并在正文口径上概括为以注意力 KL 主导的 K 路径；而 $p_V$ 则由 per-token Value 路径上的独立输出扰动代理确定。也就是说，\texttt{RoleAlign} 在低比特场景中保持统一的是流程，而不是把 K/V 两侧压成同一个单一评分函数。

在这一参数化下，
\begin{equation}
(p_K,p_V)=(100,100)
\end{equation}
对应标准 min-max 非对称量化端点。该端点解释的是格式空间本身：它说明 \texttt{RoleAlign} 的分位数参数化严格包含了最朴素的 min-max 非对称表示作为边界情形，从而使本文后续与 \texttt{KIVI-style} 的比较能够建立在同一角色感知格式族下，而不是建立在跨格式比较之上。

% 图 3-5 建议置于本小节第 1 段之后、第 2 段之前
% 图题：图 3-5 INT4-RoleAlign 的角色感知非对称量化轴
\begin{figure}[!htbp]
  \centering
  \begin{tikzpicture}[
    every node/.style={font=\small, align=center},
    mat/.style={rectangle, draw=blue!55!black, thick, minimum width=2.5cm, minimum height=2.5cm},
    note/.style={rectangle, rounded corners=3pt, draw=gray!70, thick, fill=gray!10,
                 minimum width=2.9cm, minimum height=0.9cm},
    arr/.style={-stealth, thick}
  ]
    % K side
    \node[font=\bfseries] at (-4.2,3.2) {Key 侧：优先保持排序结构};
    \node[font=\scriptsize] at (-4.2,2.65) {token 聚合 $\rightarrow$ channel 独立量化参数};
    \node[mat] (kmat) at (-4.2,1.2) {};
    \foreach \x in {-5.0,-4.5,-4.0,-3.5} {\draw[gray!55] (\x,0.0) -- (\x,2.4);}
    \foreach \y in {0.6,1.2,1.8} {\draw[gray!55] (-5.4,\y) -- (-3.0,\y);}
    \node[rotate=90, font=\scriptsize] at (-5.8,1.2) {token};
    \node[font=\scriptsize] at (-4.2,-0.45) {channel};
    \draw[arr, red!70!black] (-4.8,2.8) -- (-4.8,2.0);
    \draw[arr, red!70!black] (-4.3,2.8) -- (-4.3,2.0);
    \draw[arr, red!70!black] (-3.8,2.8) -- (-3.8,2.0);
    \node[note] (kagg) at (-4.2,-1.6) {沿 token 聚合\\按 channel 独立确定 $(s_K,z_K)$};
    \node[note] (kpk) at (-4.2,-3.0) {$p_K$ percentile\\attention KL};
    \draw[arr] (kmat.south) -- (kagg.north);
    \draw[arr] (kagg.south) -- (kpk.north);

    % V side
    \node[font=\bfseries] at (4.2,3.2) {Value 侧：优先保持内容聚合};
    \node[font=\scriptsize] at (4.2,2.65) {channel 聚合 $\rightarrow$ token 独立量化参数};
    \node[mat] (vmat) at (4.2,1.2) {};
    \foreach \x in {3.4,3.9,4.4,4.9} {\draw[gray!55] (\x,0.0) -- (\x,2.4);}
    \foreach \y in {0.6,1.2,1.8} {\draw[gray!55] (3.0,\y) -- (5.4,\y);}
    \node[rotate=90, font=\scriptsize] at (2.6,1.2) {token};
    \node[font=\scriptsize] at (4.2,-0.45) {channel};
    \draw[arr, orange!80!black] (3.2,1.8) -- (4.0,1.8);
    \draw[arr, orange!80!black] (3.2,1.2) -- (4.0,1.2);
    \draw[arr, orange!80!black] (3.2,0.6) -- (4.0,0.6);
    \node[note] (vagg) at (4.2,-1.6) {沿 channel 聚合\\按 token 独立确定 $(s_V,z_V)$};
    \node[note] (vpv) at (4.2,-3.0) {$p_V$ percentile\\输出扰动代理};
    \draw[arr] (vmat.south) -- (vagg.north);
    \draw[arr] (vagg.south) -- (vpv.north);
  \end{tikzpicture}
  \caption{INT4-RoleAlign 的角色感知非对称量化轴。Key 侧沿 token 方向聚合、对每个 channel 单独确定非对称量化参数，以优先保持 attention ranking；Value 侧沿 channel 方向聚合、对每个 token 单独确定非对称量化参数，以适应内容侧的动态范围变化。}
  \label{fig:ch3-rolealign-axis}
\end{figure}


综上，\texttt{RoleAlign} 的核心不是简单地对 $K$ 和 $V$ 采用不同公式，而是把第 \ref{sec:ch3-motivation-kv} 节识别出的角色不对称，转写为量化轴不对称：Key 侧优先保持注意力排序，Value 侧在保持可聚合性的前提下承受更激进的压缩。本文最终主线配置因此固定为：K-path 由注意力 KL 主导、V-path 由独立输出扰动代理确定，运行时采用 per-channel Key + per-token Value 非对称量化。逐头温度校正在非对称格式下未观察到稳定收益，已在第 \ref{sec:ch3-calibration} 节被保留为补充诊断对象；相关推导与溯源说明见附录第~\ref{sec:app-invtau-diagnostic}~节，主线默认不启用。

\subsection{INT4-RoleAlign 与 KIVI-style 的设计差异}
\label{subsec:ch3-rolealign-vs-kivi}

由于 \texttt{INT4-RoleAlign} 与 \texttt{KIVI-style} 共享角色差异化的量化轴，即均采用 per-channel Key 与 per-token Value 的非对称格式，因此有必要在方法章中明确二者到底共享什么、又在哪一层分叉。否则，后续第四章中的跨模型对比容易被误读为“两个完全不同的方法体系之间的竞赛”，而不是“同一角色感知格式族下，两种校准哲学的比较”。

从格式层看，二者的共同点十分明确：都承认 $K$ 与 $V$ 在低比特场景下不应被统一对待，并都采用 $K$ 侧更强调通道异质性、$V$ 侧更强调时序动态范围的角色差异化表示。真正的分叉发生在校准哲学与接口层。\texttt{KIVI-style} 的核心是直接依赖运行时统计量确定量化范围，不要求离线校准产物；而 \texttt{INT4-RoleAlign} 则把 K/V 两侧都纳入离线、可审计、可复现的角色感知校准接口：其中 Key 侧以注意力 KL 为主代理，并在离线搜索纪律下确定参数，Value 侧由独立输出扰动代理主导，并最终将两侧参数与量化轴元数据共同固化为可部署的校准产物。也正因为如此，\texttt{RoleAlign} 不只是一个“选择了不同 percentile”的技巧，而是一套可向预算分配与系统落地自然延伸的低比特路径。

进一步地，\texttt{RoleAlign} 与 \texttt{KIVI-style} 的关系不应被理解为“两个完全异构的方法竞争”，而应被理解为：在共享角色差异化量化格式的前提下，本文将校准哲学从运行时统计提升为离线角色感知校准，并由此获得可审计的产物接口与可扩展的后续统计读数。两者的差异主要不在格式层，而在于是否显式把行为引导原则写入离线校准流程。正因为这一边界足够清楚，第四章第 \ref{subsec:exp-int4-cross-model} 节中的对比结果才可以被解释为“在同一角色感知格式族下，离线行为引导校准是否带来额外价值”，而不是被误读为跨格式、跨实现、乃至跨问题设定的失配比较。

\begin{table}[!htbp]
  \centering
  \caption{\texttt{INT4-RoleAlign} 与 \texttt{KIVI-style} 的设计边界比较}
  \label{tab:ch3-rolealign-vs-kivi}
  \label{tab:s3-rolealign-vs-kivi}
  \scriptsize
  \setlength{\tabcolsep}{3.2pt}
  \renewcommand{\arraystretch}{1.10}
  \begin{tabularx}{\linewidth}{>{\raggedright\arraybackslash}p{1.35cm}>{\raggedright\arraybackslash}p{2.10cm}>{\raggedright\arraybackslash}p{2.35cm}X}
    \toprule
    维度 & \texttt{KIVI-style} & \texttt{RoleAlign} & 边界读法 \\
    \midrule
    量化轴 & \makecell[l]{per-channel K +\\per-token V} & 相同 & 二者处在同一低比特格式族下 \\
    参数来源 & 运行时统计 & \makecell[l]{离线 K/V\\分路径校准} & 比较集中在参数来源，而非格式差异 \\
    行为代理 & 无显式代理 & \makecell[l]{K: 注意力分布 KL\\V: 输出扰动代理} & 排序侧与内容侧分开解释 \\
    产物接口 & 无冻结产物 & 固定校准产物 & RoleAlign 可审计、可冻结、可复现 \\
    分配扩展 & 无共享画像 & 可输出共享读数 & 只说明可接入 allocator，不提前报告结果 \\
    \bottomrule
  \end{tabularx}
\end{table}

\section{行为敏感度驱动的层间预算分配与 AutoK}
\label{sec:ch3-allocator}

第 \ref{sec:ch3-calibration} 节已经将行为引导原则落实为可执行的离线校准流程，第 \ref{sec:ch3-paths} 节则给出了 INT8 基准路径与 \texttt{INT4-RoleAlign} 的跨位宽实例。在这一基础上，本节进一步回答一个更高层的问题：在给定位宽预算约束下，哪些层或哪些角色应被优先保护。也就是说，若校准层回答的是“如何量化”，那么分配层回答的就是“哪里优先保护”。本文将这一层称为行为引导分配，并强调其角色是统一框架的扩展层，而不是与校准并列竞争主轴的独立方法体系。

为了统一记号，设共有 $L$ 层，第 $l$ 层的 Key 与 Value 位宽分别记为 $b_K^{(l)}$ 和 $b_V^{(l)}$，则平均 KV bit 预算可写为
\begin{equation}
\bar b
=
\frac{1}{2L}\sum_{l=1}^{L}\bigl(b_K^{(l)}+b_V^{(l)}\bigr).
\end{equation}
在对称分配的特殊情形下，若 $b_K^{(l)}=b_V^{(l)}=b^{(l)}$，则有
\begin{equation}
\bar b
=
\frac{1}{L}\sum_{l=1}^{L} b^{(l)}.
\end{equation}
后续 allocator 的任务，就是在给定 $\bar b$ 的前提下，为各层或各角色生成满足预算约束的位宽决策向量。

\subsection{逐层敏感度构建与分配策略}

行为引导分配层与校准层共享同一离线来源，但其操作对象并不是在线重新计算的注意力 KL，而是由离线校准链路产出的校准产物所诱导的层级可读数。记第 $l$ 层从校准产物中读取的 K/V 角色相关缩放张量为
\begin{equation}
\sigma_K^{(l)}, \qquad \sigma_V^{(l)},
\end{equation}
它们分别表示该层在离线校准后固化下来的 Key 与 Value 侧角色尺度读数。这里需要明确两点：其一，这些读数来自离线校准产物，而不是在分配阶段重新对样本逐 token / 逐 head 计算注意力 KL；其二，它们在本节中承担的是行为敏感度的 \emph{可操作代理} 角色，而不是被宣称为与校准层的注意力 KL 数值等价。

在最基础的逐层 allocator 中，本文采用 K 侧角色尺度作为默认的逐层敏感度代理。其理由与第 \ref{sec:ch3-motivation-kv} 节中的低比特动机诊断一致：在当前研究族与本文关注的低比特场景下，Key 精度下降是导致对称 \texttt{INT4} 崩塌的主导因素，因此优先使用 K 侧读数构建逐层画像更符合当前预算分配的结构需要。这里的写法应被理解为工作性代理，而不是对所有模型家族、模型规模与 GQA 结构都成立的普适 allocator 原理。记聚合模式为 $\alpha\in\{\max,\mathrm{mean}\}$，则第 $l$ 层的逐层敏感度写为
\begin{equation}
s_{\alpha}^{(l)}
=
\operatorname{Agg}_{\alpha}\!\bigl(\sigma_K^{(l)}\bigr),
\qquad l=1,\dots,L.
\end{equation}
相应地，整层画像记为
\begin{equation}
S_{\alpha}
=
\{s_{\alpha}^{(1)},s_{\alpha}^{(2)},\dots,s_{\alpha}^{(L)}\}.
\end{equation}
这里的 $\operatorname{Agg}_{\alpha}$ 表示对除层索引外的所有轴做一次性聚合；当 $\alpha=\max$ 时，对应“最脆弱读数主导”的保守口径，当 $\alpha=\mathrm{mean}$ 时，对应“整体平均敏感度”口径。后文实验中将分别以 BA-$k_{\max}$ 与 BA-$k_{\mathrm{mean}}$ 标记这两类逐层分配器。

给定固定保护层数 $k$，基础逐层 allocator 首先定义保护层集合
\begin{equation}
\mathcal P_k^{(\alpha)}
=
\operatorname{TopK}(S_{\alpha},k),
\end{equation}
再将位宽决策写成两级向量
\begin{equation}
b^{(l)}
=
b_{\mathrm{lo}}
+
\bigl(b_{\mathrm{hi}}-b_{\mathrm{lo}}\bigr)\,
\mathbf 1\!\left[l\in \mathcal P_k^{(\alpha)}\right].
\end{equation}
此时，平均预算满足
\begin{equation}
\bar b
=
b_{\mathrm{lo}}
+
\frac{k}{L}\bigl(b_{\mathrm{hi}}-b_{\mathrm{lo}}\bigr).
\end{equation}
这说明 $k$ 在本节中不应被理解为一个任意经验数字，而应被理解为在固定两级位宽设置下控制预算强度的结构参数。BA-$k$ 的本质不是“保护前几层”，而是在预算受限时保护行为代理最高的若干层。

% 图 3-6 建议置于本小节后半段
% 图题：图 3-6 行为敏感度画像到预算分配的两级流程：逐层分配与角色扩展
\begin{figure}[!htbp]
  \centering
  \begin{tikzpicture}[
    x=1cm, y=1cm,
    >=Stealth,
    every node/.style={font=\rmfamily\small, align=center, text=black},
    arr/.style={-stealth, line width=1.05pt, draw=black!78},
    leftpanel/.style={
      rectangle, rounded corners=7pt,
      draw=black!70!cyan, line width=1.0pt,
      fill=cyan!4,
      minimum width=5.95cm, minimum height=8.30cm
    },
    rightpanel/.style={
      rectangle, rounded corners=7pt,
      draw=black!55, dashed, line width=1.0pt,
      fill=gray!3,
      minimum width=5.95cm, minimum height=8.30cm
    },
    titleL/.style={font=\bfseries\small, text=black!80!cyan},
    titleR/.style={font=\bfseries\small, text=black!72},
    subtitleL/.style={font=\rmfamily\scriptsize, text=black!65!cyan},
    subtitleR/.style={font=\rmfamily\scriptsize, text=black!58},
    tealcard/.style={
      rectangle, rounded corners=5pt,
      draw=black!70!cyan, line width=0.95pt,
      fill=white, text width=3.75cm,
      minimum height=1.02cm, inner sep=5pt
    },
    profilecard/.style={
      rectangle, rounded corners=5pt,
      draw=black!70!cyan, line width=0.95pt,
      fill=white, text width=3.75cm,
      minimum height=1.32cm, inner sep=5pt
    },
    tealslim/.style={
      rectangle, rounded corners=5pt,
      draw=black!70!cyan, line width=0.95pt,
      fill=white, text width=3.15cm,
      minimum height=0.88cm, inner sep=5pt
    },
    redcard/.style={
      rectangle, rounded corners=5pt,
      draw=red!58!black, line width=0.95pt,
      fill=white, text width=3.55cm,
      minimum height=0.94cm, inner sep=5pt
    },
    sandcard/.style={
      rectangle, rounded corners=5pt,
      draw=black!52, line width=0.9pt,
      fill=white, text width=3.55cm,
      minimum height=1.02cm, inner sep=5pt
    },
    branchL/.style={
      rectangle, rounded corners=5pt,
      draw=black!70!cyan, line width=0.95pt,
      fill=white, text width=2.25cm,
      minimum height=1.16cm, inner sep=5pt
    },
    branchR/.style={
      rectangle, rounded corners=5pt,
      draw=black!52, line width=0.9pt,
      fill=white, text width=2.25cm,
      minimum height=1.16cm, inner sep=5pt
    },
    tiny/.style={font=\rmfamily\scriptsize}
  ]

    \node[font=\bfseries\large] at (0,5.02)
      {行为敏感度画像到预算决策};
    \node[font=\rmfamily\small, text=black!62] at (0,4.52)
      {逐层主线与 K/V 角色预算接口};
    \draw[brown!55!black, line width=0.8pt] (-5.95,4.10) -- (5.95,4.10);

    \node[leftpanel] (lp) at (-3.45,0.12) {};
    \node[rightpanel] (rp) at (3.45,0.12) {};

    % Left title
    \node[titleL] at (-3.45,3.28) {逐层分配};
    \node[subtitleL] at (-3.45,2.88) {从 readout 走到 layer-wise budget};

    % Left vertical chain
    \node[tealcard] (readout) at (-3.45,1.80)
      {artifact-based readout\\[2pt]$\sigma_K^{(l)},\ \sigma_V^{(l)}$};

    \node[tealslim] (agg) at (-3.45,0.68)
      {$\operatorname{Agg}_{\alpha}$，\ $\alpha\in\{\max,\mathrm{mean}\}$};

    \node[profilecard] (profile) at (-3.45,-0.63)
      {layer sensitivity profile\\[2pt]$S=\{s^{(l)}\}$};

    % simple bars inside profile card
    \foreach \x/\h in {-4.52/0.22,-4.31/0.62,-4.10/0.40,-3.89/0.82,-3.68/0.48,-3.47/0.30,-3.26/0.76,-3.05/0.45,-2.84/0.58} {
      \fill[cyan!18] (\x,-1.48) rectangle ++(0.10,\h);
      \draw[black!45!cyan, line width=0.32pt] (\x,-1.48) rectangle ++(0.10,\h);
    }

    \node[tealslim] (topk) at (-3.45,-2.22)
      {Top-$k$ 保护集合\\[1pt]$\mathcal P_k=\operatorname{TopK}(S,k)$};

    \node[redcard] (bits) at (-3.45,-3.48)
      {layer-wise bit vector\\[2pt]$[b^{(1)},\dots,b^{(L)}]$};

    \draw[arr] (readout.south) -- (agg.north);
    \draw[arr] (agg.south) -- (profile.north);
    \draw[arr] (profile.south) -- (topk.north);
    \draw[arr] (topk.south) -- (bits.north);

    % Right title
    \node[titleR] at (3.45,3.28) {接口扩展};
    \node[subtitleR] at (3.45,2.88) {只定义 K/V 角色预算接口，不作为主实证策略};

    \node[sandcard] (shared) at (3.45,1.86)
      {可选角色条件化 readout\\[2pt]$(s_K^{(l)},\ s_V^{(l)})$};

    \node[branchL] (sk) at (2.08,-0.08)
      {$S_K$\\[2pt]Key budget\\interface};
    \node[branchR] (sv) at (4.82,-0.08)
      {$S_V$\\[2pt]Value budget\\interface};

    \node[branchL] (pk) at (2.08,-1.92)
      {$\mathcal P_K$\\[2pt]Key mask};
    \node[branchR] (pv) at (4.82,-1.92)
      {$\mathcal P_V$\\[2pt]Value mask};

    \node[redcard] (pair) at (3.45,-3.58)
      {pair-wise budget\\interface\\[2pt]$\mathbf b^{(l)}=(b_K^{(l)},\,b_V^{(l)})$};

    \coordinate (split) at (3.45,0.82);
    \coordinate (merge) at (3.45,-2.78);

    \draw[arr] (shared.south) -- (split);
    \draw[arr] (split) -- (sk.north);
    \draw[arr] (split) -- (sv.north);
    \draw[arr] (sk.south) -- (pk.north);
    \draw[arr] (sv.south) -- (pv.north);
    \draw[arr] (pk.south) -- (merge);
    \draw[arr] (pv.south) -- (merge);
    \draw[arr] (merge) -- (pair.north);

  \end{tikzpicture}
  \caption{行为敏感度画像到预算分配的两级流程：逐层主线与角色预算接口。左半部分给出本文主线实证使用的逐层分配链路，从 artifact-based readout 到排序、top-$k$ 保护集合与 layer-wise bit vector；右半部分只说明相同框架可以扩展到 K/V 角色预算接口，不构成本文已经完整验证的独立 role-aware allocator。}
  \label{fig:ch3-allocation-flow}
\end{figure}


\subsection{K/V 角色感知的预算扩展}

若进一步考虑到第 \ref{sec:ch3-motivation-kv} 节所揭示的 K/V 不对称性，则逐层预算还可以自然扩展为角色感知预算。此时，第 $l$ 层不再仅对应一个标量位宽，而是对应一个角色对
\begin{equation}
\mathbf b^{(l)}
=
\bigl(b_K^{(l)},\, b_V^{(l)}\bigr).
\end{equation}
与前述对称分配情形一致，全局平均预算仍满足
\begin{equation}
\bar b
=
\frac{1}{2L}\sum_{l=1}^{L}
\left(
b_K^{(l)}+b_V^{(l)}
\right).
\end{equation}

为了表达角色感知扩展，本文进一步引入 K/V 两侧的角色条件化敏感度读数：
\begin{equation}
s_K^{(l)}, \qquad s_V^{(l)},
\end{equation}
它们表示第 $l$ 层在 K 与 V 两侧的角色条件化敏感度读数，可由同一离线校准来源中的角色条件化 traces 或诊断读数实例化，而不要求在本小节将其唯一地锁定为某一种估计量。于是，角色感知分配器可抽象记为
\begin{equation}
\mathbf b^{(l)}
=
\mathcal A_{KV}
\bigl(
s_K^{(l)},\,s_V^{(l)};\bar b
\bigr).
\end{equation}
在当前实现中，这类读数可以由基于校准产物的角色尺度聚合给出；但在本节的框架层表达中，更重要的是说明行为引导分配可以自然扩展到 K/V 角色条件化的预算接口，而不是把某一种具体读数形式误写成唯一主定义。

在最简单的两级版本中，可分别定义 K 与 V 的保护集合：
\begin{equation}
\mathcal P_K^{(\alpha)}
=
\operatorname{TopK}(S_{K,\alpha},k_K),
\qquad
\mathcal P_V^{(\alpha)}
=
\operatorname{TopK}(S_{V,\alpha},k_V),
\end{equation}
并写成
\begin{equation}
b_K^{(l)}
=
b_{K,\mathrm{lo}}
+
\bigl(b_{K,\mathrm{hi}}-b_{K,\mathrm{lo}}\bigr)
\mathbf 1\!\left[l\in\mathcal P_K^{(\alpha)}\right],
\end{equation}
\begin{equation}
b_V^{(l)}
=
b_{V,\mathrm{lo}}
+
\bigl(b_{V,\mathrm{hi}}-b_{V,\mathrm{lo}}\bigr)
\mathbf 1\!\left[l\in\mathcal P_V^{(\alpha)}\right].
\end{equation}

需要强调的是，这一角色感知预算扩展在本文中主要承担“接口扩展”的角色：它说明行为引导框架并不只允许逐层一维预算，也允许沿 K/V 角色再细化一步。但在当前论文主线中，真正承担主要实证任务的仍是逐层分配器；角色感知预算在这里被写成统一框架向更细粒度预算接口的自然生长，而不是一个已经被完整证明为普遍更优的独立主张。

\subsection{敏感度聚合、鲁棒性与覆盖度准则}

前两小节已经给出了由校准产物到逐层/角色感知预算的基本接口，但原始角色尺度读数仍需要经过聚合与归一化，才能被转写为稳定的预算表达。本文在实现层面采用的聚合模式并不是“token 先均值、head 再取最大、sample 再取均值”这样的嵌套规则，而是对除 layer 外的所有轴一次性聚合。对任意层级读数 $\sigma^{(l)}$，定义
\begin{equation}
\operatorname{Agg}_{\alpha}\!\bigl(\sigma^{(l)}\bigr)
=
\begin{cases}
\displaystyle \max_{u\in \mathcal U_l}\sigma^{(l)}_u, & \alpha=\max,\\[6pt]
\displaystyle \frac{1}{|\mathcal U_l|}\sum_{u\in\mathcal U_l}\sigma^{(l)}_u, & \alpha=\mathrm{mean},
\end{cases}
\end{equation}
其中 $\mathcal U_l$ 表示第 $l$ 层除层索引外的所有索引集合。$\alpha=\max$ 与 $\alpha=\mathrm{mean}$ 因而应被理解为两个可选的聚合模式，而不是彼此嵌套的三层操作。后续实验也将分别以 BA-$k_{\max}$ 与 BA-$k_{\mathrm{mean}}$ 对应这两类配置。

为比较不同模型的 profile 形状，本文进一步定义归一化敏感度
\begin{equation}
\tilde s^{(l)}
=
\frac{\max(s^{(l)},0)}
{\sum_{j=1}^{L}\max(s^{(j)},0)},
\qquad
\sum_{l=1}^{L}\tilde s^{(l)}=1.
\end{equation}
将其按降序排列为
\begin{equation}
\tilde s_{(1)}\ge \tilde s_{(2)}\ge \cdots \ge \tilde s_{(L)},
\end{equation}
即可得到累计覆盖度函数
\begin{equation}
\Gamma(k)
=
\sum_{i=1}^{k}\tilde s_{(i)}.
\end{equation}
这一函数把“画像是集中还是弥散”转写成了一个预算上可解释的对象：若某模型的敏感度主要集中在少数层，则 $\Gamma(k)$ 会在较小的 $k$ 上迅速上升；若某模型的敏感度在多层之间更加弥散，则达到同一覆盖阈值所需的 $k$ 会显著增大。也就是说，覆盖度曲线不是另一个结果图，而是把敏感度画像的形状转化为预算决策表达的中间桥梁。

% 图 3-7 建议置于本小节中部
% 图题：图 3-7 敏感度覆盖度曲线示意：集中型画像与弥散型画像的预算含义
\begin{figure}[!htbp]
  \centering
  \begin{tikzpicture}[x=1.18cm,y=4.85cm,>=Stealth]
    % axes
    \draw[->, line width=1.0pt, draw=black!80] (0,0) -- (6.35,0)
      node[below, font=\small] {$k$};
    \draw[->, line width=1.0pt, draw=black!80] (0,0) -- (0,1.06)
      node[left, font=\small] {$\Gamma(k)$};

    % ticks
    \foreach \y/\lab in {0.0/0,0.4/0.4,0.8/0.8,1.0/1.0} {
      \draw[black!75] (-0.04,\y) -- (0.04,\y);
      \node[left, font=\scriptsize, text=black!72] at (-0.08,\y) {\lab};
    }
    \foreach \x in {1,2,3,4,5,6} {
      \draw[black!75] (\x,0.0) -- (\x,-0.018);
      \node[below, font=\scriptsize, text=black!72] at (\x,-0.02) {\x};
    }

    % threshold lines
    \draw[densely dashed, gray!35] (0,0.80) -- (6.0,0.80);
    \draw[densely dashed, gray!55] (0,0.85) -- (6.0,0.85);
    \draw[densely dashed, gray!35] (0,0.90) -- (6.0,0.90);
    \node[right, font=\scriptsize, text=gray!70!black] at (6.03,0.80) {$\rho=0.80$};
    \node[right, font=\scriptsize, text=gray!78!black] at (6.03,0.85) {$\rho=0.85$};
    \node[right, font=\scriptsize, text=gray!70!black] at (6.03,0.90) {$\rho=0.90$};

    % filled areas
    \fill[cyan!10]
      (0.55,0.00) --
      plot[smooth] coordinates {
        (0.55,0.24) (1.00,0.50) (1.55,0.76) (2.05,0.86)
        (2.85,0.93) (3.90,0.97) (5.60,1.00)
      } --
      (5.60,0.00) -- cycle;

    \fill[orange!10]
      (0.55,0.00) --
      plot[smooth] coordinates {
        (0.55,0.10) (1.05,0.22) (1.75,0.37) (2.55,0.52)
        (3.35,0.67) (4.15,0.79) (4.85,0.86) (5.55,0.92)
      } --
      (5.55,0.00) -- cycle;

    % curves
    \draw[cyan!65!black, line width=1.7pt]
      plot[smooth] coordinates {
        (0.55,0.24) (1.00,0.50) (1.55,0.76) (2.05,0.86)
        (2.85,0.93) (3.90,0.97) (5.60,1.00)
      };

    \draw[orange!82!black, line width=1.7pt]
      plot[smooth] coordinates {
        (0.55,0.10) (1.05,0.22) (1.75,0.37) (2.55,0.52)
        (3.35,0.67) (4.15,0.79) (4.85,0.86) (5.55,0.92)
      };

    % highlighted target and k* markers
    \fill[cyan!70!black] (2.05,0.85) circle (1.6pt);
    \fill[orange!82!black] (4.85,0.85) circle (1.6pt);
    \draw[densely dotted, cyan!65!black, line width=0.9pt] (2.05,0.00) -- (2.05,0.85);
    \draw[densely dotted, orange!82!black, line width=0.9pt] (4.85,0.00) -- (4.85,0.85);
    \fill[cyan!65!black] (2.05,0.00) circle (1.35pt);
    \fill[orange!82!black] (4.85,0.00) circle (1.35pt);

    % labels
    \node[font=\scriptsize, text=cyan!70!black,
          fill=white, inner sep=2pt, rounded corners=3pt]
      at (1.95,0.97) {集中型画像};
    \node[font=\scriptsize, text=orange!82!black,
          fill=white, inner sep=2pt, rounded corners=3pt]
      at (5.02,0.60) {弥散型画像};

    \node[font=\scriptsize, text=gray!75!black,
          fill=white, inner sep=2pt]
      at (5.08,0.875) {相同 coverage 目标};

    \node[font=\scriptsize, text=cyan!70!black] at (2.05,-0.08) {$k^\star_{\text{集中}}$};
    \node[font=\scriptsize, text=orange!82!black] at (4.85,-0.08) {$k^\star_{\text{弥散}}$};
  \end{tikzpicture}
  \caption{敏感度覆盖度曲线示意：集中型画像与弥散型画像的预算含义。相同覆盖阈值下，集中型画像会在更小的 $k$ 上达到目标覆盖度，而弥散型画像需要更大的保护层数；这正是 AutoK 将 profile 形状转写为预算建议的核心直觉。}
  \label{fig:ch3-coverage-curve}
\end{figure}


从方法层的角度看，$\Gamma(k)$ 的价值在于：它既保留了逐层敏感度排序，又避免将预算建议器写成“固定保护前 $k$ 层”的经验规则。正因为如此，本文后续才可以把 \texttt{AutoK} 写成敏感度画像的自然输出，而不是另起炉灶再设计一套与前文脱节的预算调节器。

\subsection{AutoK 的预算自动建议机制}
\label{sec:ch3-autok}

在固定 $k$ 的 BA-$k$ 中，预算强度需要由人工扫描决定，而不同模型的敏感度画像形状并不一致，这意味着相同的 $k$ 未必对应相同程度的“行为覆盖”。为减少这种人工扫描负担，本文引入 \texttt{AutoK} 作为由画像驱动的预算建议器：它不改变第 \ref{sec:ch3-calibration} 节的校准原则，也不重新训练任何额外模块，而只是将前述覆盖度准则直接转写为最小保护层数建议。

给定覆盖阈值 $\rho\in(0,1)$，定义
\begin{equation}
k^\star(\rho)
=
\min\left\{
k:\Gamma(k)\ge \rho
\right\}.
\end{equation}
于是，\texttt{AutoK} 对应的保护集合写为
\begin{equation}
\mathcal P_{\mathrm{AutoK}}^{(\alpha,\rho)}
=
\operatorname{TopK}\!\bigl(S_{\alpha},\,k^\star(\rho)\bigr),
\end{equation}
对应的两级位宽向量为
\begin{equation}
b_{\mathrm{AutoK}}^{(l)}
=
b_{\mathrm{lo}}
+
\bigl(b_{\mathrm{hi}}-b_{\mathrm{lo}}\bigr)
\mathbf 1\!\left[l\in \mathcal P_{\mathrm{AutoK}}^{(\alpha,\rho)}\right].
\end{equation}
在图示和后续实验中，覆盖度预算建议统一以阈值 $\rho$ 表示；例如 $\rho=0.80$、$0.85$ 与 $0.90$ 分别表示覆盖累计敏感度质量的 80\%、85\% 与 90\%。

需要明确的是，\texttt{AutoK} 的定位是预算建议器，而不是“第三个主方法”或“对所有模型稳定取胜的通用策略”。它的价值在于：将敏感度画像的形状直接转写为无需人工扫 $k$ 的预算建议，从而使同一行为引导框架在跨模型迁移时具备更强的可读性与可复用性。至于不同聚合模式、不同覆盖阈值是否在所有模型家族、模型规模与任务上都表现一致，则属于后续第四章的实证问题，而不在本节作先验结论。

% 表 3-4 建议置于本小节末尾
% 表题：表 3-4 分配层策略的输入、选择规则与对应记号
\begin{table}[!htbp]
  \centering
  \caption{分配层策略的输入、选择规则与对应记号}
  \label{tab:ch3-allocator-policies}
  \footnotesize
  \setlength{\tabcolsep}{3.6pt}
  \renewcommand{\arraystretch}{1.10}
  \begin{tabularx}{\linewidth}{>{\raggedright\arraybackslash}p{1.35cm}>{\raggedright\arraybackslash}p{1.3cm}X>{\centering\arraybackslash}p{0.95cm}>{\raggedright\arraybackslash}p{1.35cm}>{\raggedright\arraybackslash}p{1.6cm}X}
    \toprule
    \textbf{策略} & \textbf{输入} & \textbf{保护规则} & \textbf{扫描?} & \textbf{输出} & \textbf{实验记号} & \textbf{方法定位} \\
    \midrule
    Uniform & 无 & 所有层统一默认低 bit & 否 & 统一预算向量 & Uniform & 低预算参考 \\
    BA-$k$ & $(\mathcal S, k)$ & 保护 sensitivity 最高的 $k$ 层 & 是 & 分层 bit 向量 & BA-k & 分层 allocator 主体 \\
    Heuristic-$k$ & $(L, k)$ & 按固定位置规则保护 $k$ 层 & 是 & 分层 bit 向量 & Heuristic-k & 强基线 \\
    BA-AutoK & $(\mathcal S, \rho)$ & 取满足 coverage 的最小 $k^\star$ & 否 & 分层 bit 向量 & BA-AutoK & coverage 预算建议 \\
    \bottomrule
  \end{tabularx}
\end{table}

\section{系统级部署实现与开销分析}
\label{sec:ch3-deployment}
\label{sec:ch3-triton}

前述各节已经分别给出了行为引导校准目标、跨位宽路径实例以及预算分配接口。本节进一步回答一个实现层问题：这些离线定义好的量化规则与预算决策，如何被组织为可部署、可复现、可审计的推理系统语义。本节的范围是离线校准产物、在线推理管线、融合解码核函数的数据流，以及相应的理论开销；所有关于 TPOT、性能交叉边界、不同模型之间速度优劣与部署建议的实测结果，统一留到第四章的部署效率小节讨论。

\subsection{离线校准产物与在线推理管线}

行为引导框架采用“离线校准、在线执行”的两阶段分工。离线阶段利用少量校准样本执行第 \ref{sec:ch3-calibration} 节定义的搜索与选择纪律，生成路径特定的校准产物；在线阶段只读取这些固定产物并执行既定量化规则，而不再重新组织目标函数，也不在推理关键路径中重复进行参数搜索。这样的分工使得行为引导原则既能保留离线校准的可审计性，也不会把额外的优化开销转移到在线解码阶段。

为统一表述，记路径相关的离线校准产物为
\begin{equation}
\mathcal A_{\mathrm{path}}
=
\left\{
\theta_{\mathrm{path}}^{(l)},
m_{\mathrm{path}}^{(l)}
\right\}_{l=1}^{L},
\end{equation}
其中 $\theta_{\mathrm{path}}^{(l)}$ 表示第 $l$ 层的路径相关量化参数，$m_{\mathrm{path}}^{(l)}$ 表示其元数据，例如量化轴、分组大小、保护裕度与模式标识等。对 INT8 基准路径而言，产物主要包括逐层逐组的静态缩放参数、保护参数与分组元数据；对对称 \texttt{INT4} 试探路径而言，其接口与 INT8 基准路径相似，但载荷与量化网格更低比特化；对 \texttt{INT4-RoleAlign} 而言，产物则组织为角色感知的校准产物，包括 K/V 的分位数参数、量化轴元数据，以及补充诊断所需的元数据。需要强调的是，这些校准产物在运行时是只读对象，其作用是冻结离线阶段已经确定的行为引导决策，而不是在推理时继续自适应搜索。

在线推理阶段的核心语义是“写入即冻结”。设第 $t$ 个新 token 在第 $l$ 层待写入缓存的张量为 $x_t^{(l)}$，则当前步的量化写入可写为
\begin{equation}
\hat x_t^{(l)}
=
Q_{\theta_t^{(l)}}\!\left(x_t^{(l)}\right),
\end{equation}
其中 $\theta_t^{(l)}$ 由离线校准产物与当前运行时上下文共同决定。这里需要特别说明两点：其一，历史 token 一旦以某组量化参数写入缓存，其量化值与对应元数据不再被回写；其二，自适应保护若在当前时刻触发，只影响当前 token 的写入，不会逆向修改已存历史缓存。这一“append 时确定、历史不回写”的语义，是本文后续所有缓存管理与融合核函数设计的基础。

从系统流程上看，Prefill 阶段根据路径规则将当前块的 K/V 写入量化缓存，Decode 阶段则读取历史量化缓存，并将其路由到对应的注意力后端。由于在线阶段不再重跑离线校准，推理系统真正需要解决的问题就不再是“如何选参数”，而是“如何让低比特缓存尽可能直接地参与解码注意力计算”。这正是下一小节 Triton 融合解码核函数所要承担的角色。

若把上述部署闭环改写为系统语义，则其核心链路可以概括为：离线阶段基于校准样本执行行为引导搜索，并将结果固化为只读的路径特定产物 $\mathcal A_{\mathrm{path}}$；在线阶段的 Prefill 只负责依据该产物完成量化写入，Decode 则只负责读取历史量化缓存并路由到对应后端执行。运行时真正被保留下来的不是离线搜索过程本身，而是其冻结后的参数接口、元数据接口以及“append 时确定、历史不回写”的缓存语义。正因为这一闭环已经在系统层被固定下来，下一小节才可以把关注点进一步收缩到：这些低比特缓存如何在 Decode 阶段以尽量少的中间物化直接参与注意力计算。

% 表 3-5 建议置于本小节末尾
% 表题：表 3-5 不同量化路径的离线产物、运行时缓存载荷与默认执行后端
\begin{table}[!htbp]
  \centering
  \caption{不同量化路径的离线产物、运行时缓存载荷与默认执行后端}
  \label{tab:ch3-runtime-paths}
  \scriptsize
  \setlength{\tabcolsep}{2.0pt}
  \renewcommand{\arraystretch}{1.14}
  \begin{tabularx}{\linewidth}{>{\raggedright\arraybackslash}p{1.35cm}>{\raggedright\arraybackslash}p{2.35cm}>{\raggedright\arraybackslash}p{2.95cm}>{\raggedright\arraybackslash}p{2.25cm}X}
    \toprule
    路径 & 离线产物 & 追加写入 & 解码后端 & 边界说明 \\
    \midrule
    INT8 & scale、保护裕度、轴元数据 & INT8 K/V 与 scale & Triton 融合后端 & 主部署路径 \\
    对称 INT4 & scale、轴元数据 & packed INT4 K/V & 参考实现或 Triton 试探 & 压力测试；解包计入成本 \\
    RoleAlign & $(p_K,p_V)$、K/V 轴、头数 & K per-channel；V per-token & 参考实现；融合扩展 & 质量核对优先；速度见第四章 \\
    \bottomrule
  \end{tabularx}
\end{table}


\subsection{Triton 融合解码核函数设计}
\label{subsec:ch3-triton-int4-asym}

在自回归解码阶段，单步计算通常只涉及一个新的 query token，但需要访问全部历史 KV Cache，因此其瓶颈主要在访存而非算术量。若采用朴素路径，系统需要先从低比特缓存中读出 K/V，将其完整反量化为中间 FP16 张量，再把这些中间张量送入标准注意力计算。该过程虽然在功能上是正确的，但会引入额外的中间张量物化与高带宽显存读写开销，从而放大 decode 阶段的 memory-bound 特征。本文在 INT8 基准路径上采用 Triton 融合解码核函数，目的正是在不改变注意力算法本身的前提下，尽量消除这类不必要的中间物化。

对 \texttt{INT8} 路径而言，融合核函数的主数据流可以概括为以下四步。首先，从量化缓存中按 block 读取一块低比特 $K_B, V_B$ 及其对应的缩放参数与元数据；其次，在片上存储（SRAM / register）内完成解包与反量化；随后，直接在片上计算当前 query 与该块 Key 的点积；最后，利用在线 softmax 维护跨块递推状态，并将结果直接累加到输出而不显式物化完整注意力矩阵。设第 $r$ 个 block 的在线 softmax 状态分别记为最大值、归一化和输出累加器
\[
m^{(r)}, \qquad \ell^{(r)}, \qquad o^{(r)},
\]
则跨块更新可写为
\begin{equation}
m^{(r+1)}
=
\max\!\left(m^{(r)},\, m_B\right),
\end{equation}
\begin{equation}
\ell^{(r+1)}
=
e^{m^{(r)}-m^{(r+1)}}\ell^{(r)}
+
\sum_{i\in B} e^{z_i-m^{(r+1)}},
\end{equation}
\begin{equation}
o^{(r+1)}
=
\frac{
e^{m^{(r)}-m^{(r+1)}}\ell^{(r)}o^{(r)}
+
\sum_{i\in B} e^{z_i-m^{(r+1)}} v_i
}{
\ell^{(r+1)}
}.
\end{equation}
因此，融合的关键收益不在于改变 softmax 的定义，而在于把“反量化 $\rightarrow$ 点积 $\rightarrow$ online softmax $\rightarrow$ 输出累加”压进同一片上流水线中，从而减少全局内存传输。

对 \texttt{INT4} 路径而言，注意力算法本身并未发生变化，新增复杂度主要来自数据表示。具体地，\texttt{INT4} 采用 nibble packing，即两个 4-bit 整数打包到一个 8-bit 字节中；因此核函数的前端需要先完成半字节解包与符号扩展，随后再复用与 \texttt{INT8} 类似的“片上反量化 $\rightarrow$ 点积 $\rightarrow$ online softmax $\rightarrow$ 累加输出”主流水线。从系统视角看，这意味着 \texttt{INT4} 的额外复杂度主要是数据布局与解包开销，而不是注意力计算逻辑的本质改变。正文在此只描述当前默认实现路径；关于融合核函数的不同迭代版本、命名与实现差异，统一保留到附录第~\ref{sec:app-triton-variants}~节说明。

为了兼容 GQA，融合核函数还需要将 query heads 映射到共享的 KV heads。设 query 头数为 $H_q$、KV 头数为 $H_{kv}$，则重复因子为
\begin{equation}
N_{\mathrm{rep}}=\frac{H_q}{H_{kv}},
\end{equation}
对应的头映射写为
\begin{equation}
h_{kv}
=
\left\lfloor
\frac{h_q}{N_{\mathrm{rep}}}
\right\rfloor.
\end{equation}
这样，核函数可以按 $(B,H_q)$ 或其等价 grid 组织并行，而无需显式复制 KV heads。该映射的作用在于在保持 GQA 结构语义正确的同时，使 query heads 与共享 KV heads 之间的访问关系在融合路径中得到直接表达。

最后需要再次诚实说明的是：本文的主系统落地路径是 \texttt{INT8} 融合核函数。对 \texttt{INT4-RoleAlign} 而言，质量评测采用参考解码实现以保证语义核对；与之并行，融合后端与外部推理库适配构成可用的快速路径与实现扩展。第三章在这里只说明这些后端的系统角色，而不把其中任何一条直接写成第四章速度结论的默认来源。

若把朴素路径与融合路径并列成系统数据流，朴素路径遵循“低比特缓存 $\rightarrow$ 完整 FP16 K/V 物化 $\rightarrow$ 标准 attention”的执行顺序，融合路径则遵循“低比特 K/V 与元数据 $\rightarrow$ block load $\rightarrow$ 片上解包与反量化 $\rightarrow$ $qK^\top$ 与 online softmax $\rightarrow$ 输出累加”的执行顺序。二者的本质差异并不在 softmax 定义本身，而在于融合解码将完整 FP16 K/V 与完整注意力矩阵的物化都排除在主路径之外；与此同时，GQA 的头映射也被直接并入同一 Decode 数据流，而不是在核外先做额外重排。因此，Triton 融合核在系统层承担的是“消除中间物化、压缩全局内存传输”的角色，而不是引入另一套不同的注意力算法。

\subsection{复杂度、访存与存储开销分析}

在离线阶段，行为引导校准的主要额外代价来自候选参数扫描。对任一路径，若候选参数集合大小记为 $|\Theta_{\mathrm{path}}|$，校准样本数为 $N$，模型层数为 $L$，Query 头数为 $H_q$，每个样本的校准长度为 $n$，头维度为 $d_k$，则离线搜索复杂度可统一写为
\begin{equation}
T_{\mathrm{calib}}^{\mathrm{path}}
=
\mathcal O\!\left(
|\Theta_{\mathrm{path}}|\,N\,L\,H_q\,n\,d_k
\right).
\end{equation}
这一表达式的意义在于：不同路径的主要差异并不体现在“是否需要离线校准”，而体现在候选参数集合的大小与结构上。INT8 基准路径、对称 \texttt{INT4} 与 \texttt{INT4-RoleAlign} 共享的是离线搜索框架，而不是相同规模的搜索空间。

在缓存显存上，若只考虑基准的对称 per-group 表示，则 KV Cache 的精确存储量可写为如下三种情况。FP16 参考开销为
\begin{equation}
M_{\mathrm{FP16}}
=
4L H_{kv} S d_k,
\end{equation}
其中因子 4 来自 K/V 双矩阵与每个元素 2 字节。对称 \texttt{INT8} 路径下，若 group size 为 $g$，则缓存显存为
\begin{equation}
M_{\mathrm{INT8}}
=
2L H_{kv} S d_k
\left(
1+\frac{4}{g}
\right),
\end{equation}
其中第一项来自 \texttt{INT8} 数据载荷，第二项来自逐组缩放参数元数据。类似地，对称 \texttt{INT4} 路径下有
\begin{equation}
M_{\mathrm{INT4}}
=
2L H_{kv} S d_k
\left(
\frac{1}{2}+\frac{4}{g}
\right).
\end{equation}
这里需要说明：\texttt{RoleAlign} 的角色感知表示遵循相同的“数据载荷 + 元数据载荷”分解原则，但其精确元数据项依赖 K/V 角色相关实现接口，因此正文在本节以统一分解解释为主，而将具体实测开销统一留到第四章讨论。

从 decode 阶段的访存角度看，融合核函数的潜在收益可以通过 block 级字节量近似解释。若 block 大小记为 $B_s$，则对称 \texttt{INT8} 路径每个 block 的访存量近似为
\begin{equation}
\mathrm{Bytes}_{\mathrm{INT8}}
\approx
2B_s d_k
\left(
1+\frac{4}{g}
\right),
\end{equation}
而 packed \texttt{INT4} 路径近似为
\begin{equation}
\mathrm{Bytes}_{\mathrm{INT4}}
\approx
2B_s d_k
\left(
\frac{1}{2}+\frac{4}{g}
\right).
\end{equation}
若块级浮点计算量近似按 $4B_s d_k$ 估计，则可得到相应的算术强度
\begin{equation}
AI_{\mathrm{INT8}}
\approx
\frac{2}{1+4/g},
\qquad
AI_{\mathrm{INT4}}
\approx
\frac{2}{1/2+4/g}.
\end{equation}
这些表达式本身并不给出任何实测速率结论，但它们清楚说明：在相似计算量下，低比特融合路径通过减少全局内存传输而改变了 decode 阶段的开销组成。这一分析在本节中的作用，是为系统实现的 memory-bound 特征提供理论解释，而不是在方法章内提前给出经验速度判断。

\label{subsec:ch3-system-overhead}
最后，离线校准产物本身的存储开销相对模型权重与运行时 KV Cache 而言通常可以忽略。对称 INT8 基准路径的产物规模大致为
\begin{equation}
\mathcal O\!\left(\frac{L d_k}{g}\right),
\end{equation}
而 \texttt{RoleAlign} 的产物主要由低维分位数参数与量化轴元数据构成，其量级同样远低于运行时缓存。综上，本节给出的复杂度、显存与访存公式并不承担经验结论职责，而是用于说明：行为引导框架在系统层增加的是可控的元数据与离线搜索成本，而不是一个与模型主体同量级的额外负担。更具体的部署效率读数与边界分析，将在第四章统一报告。

\section{本章小结}
\label{sec:ch3-summary}

本章围绕注意力行为保持这一统一原则，完成了从理论起点到方法框架的系统展开。首先，第 \ref{sec:ch3-problem} 节从注意力近似误差的代数分解出发，说明 KV Cache 量化真正损伤的不是孤立的张量数值距离，而是由注意力分布与聚合输出共同构成的注意力行为；第 \ref{sec:ch3-motivation-kv} 节进一步给出低比特场景下的压缩版动机诊断，指出导致对称 \texttt{INT4} 路径失稳的主导因素来自 Key 精度下降，从而为后续角色感知的低比特实例化提供了结构性依据。基于这一理论与动机基础，第 \ref{sec:ch3-framework} 节将校准与分配统一组织到同一行为引导框架下，明确了校准层回答“如何量化”，分配层回答“哪里优先保护”的两层决策关系。

在这一统一框架下，本章进一步将联合行为目标操作化为可执行的离线校准流程。第 \ref{sec:ch3-calibration} 节以注意力分布 KL 散度作为校准层的可操作代理，并给出统一的参数搜索与稳健选择纪律；第 \ref{sec:ch3-paths} 节则沿着由保守到激进的顺序，分别给出 INT8 基准路径、对称 \texttt{INT4} 的局限与格式升级动因，以及 \texttt{INT4-RoleAlign} 的角色感知非对称量化实例，同时在方法层面明确了其与 \texttt{KIVI-style} 的设计边界。进一步地，第 \ref{sec:ch3-allocator} 节将同源的行为敏感度画像延伸为逐层预算分配与 \texttt{AutoK} 预算自动建议机制，使行为引导原则从参数选择自然扩展到预算决策。

最后，第 \ref{sec:ch3-deployment} 节说明了上述方法如何被落成为可部署、可复现、可审计的系统接口，包括离线校准产物、在线推理管线、Triton 融合解码核函数以及复杂度、访存与存储开销分析。至此，本文已经在方法层回答了两个核心问题：为何应以注意力行为作为 KV Cache 量化的统一组织对象，以及这一原则如何被实例化为可执行的校准路径、预算机制与系统落地方案。下一章将在统一实验协议下，进一步验证本章提出的 INT8 基准路径保真度、低比特恢复能力、跨模型策略适用区间结构及部署边界。
